\Needspace{12\baselineskip}
\section{Proofs}
\label{app:rewrite-proofs}

\Needspace{7\baselineskip}
\begin{lemma}[Static AI-service choice]
\label{lem:rewrite-monopoly-choice}
For $K,A,L,B>0$, $0<\alpha<1$, $\sigma>0$, and strictly positive CES weights
that sum to one, problem~\eqref{eq:rewrite-service-choice} has a unique,
strictly positive global maximizer. It is characterized by
condition~\eqref{eq:rewrite-monopoly-foc} and depends smoothly on the positive
variables $K,A,L,B$, for each fixed $\sigma$.
\end{lemma}

\begin{proof}[Proof of Lemma~\ref{lem:rewrite-monopoly-choice}]
Fix \(K,A,L,B>0\) and let \(R(X)=p_X(X)X\).
The demand elasticity in \eqref{eq:rewrite-inverse-demand-elasticity} and the
CES-share derivative imply
\[
 R'(X)=p_X(X)(1-e_X),\qquad
 \frac{\dd s_X}{\dd\ln X}=\left(1-\frac1\sigma\right)s_X(1-s_X).
\]
Wherever marginal revenue is positive, direct differentiation gives
\begin{equation}
 -\frac{XR''(X)}{p_X(X)}
 =e_X(1-e_X)+\left(\alpha-\frac1\sigma\right)
 \left(1-\frac1\sigma\right)s_X(1-s_X).
 \label{eq:rewrite-monopoly-proof-curvature}
\end{equation}
If \((\alpha-1/\sigma)(1-1/\sigma)\geq0\), the right-hand side is positive.
If this product is negative, the parameter restrictions imply
\(0<\alpha<1/\sigma<1\), and
\[
 \begin{aligned}
 -\frac{XR''(X)}{p_X(X)}
 ={}&\frac1\sigma\left(1-\frac1\sigma\right)(1-s_X)^2\\
 &+\left[\frac1{\sigma^2}+2\alpha-\frac{3\alpha}{\sigma}\right]s_X(1-s_X)
 +\alpha(1-\alpha)s_X^2.
 \end{aligned}
\]
The middle coefficient is positive because
\[
 \frac1{\sigma^2}+2\alpha-\frac{3\alpha}{\sigma}
 =\left(\frac1\sigma-\alpha\right)^2
 +\alpha\left(2-\frac1\sigma-\alpha\right)>0.
\]
Thus marginal revenue decreases strictly whenever it is positive.

As \(X\downarrow0\), revenue is proportional to \(X^{1-1/\sigma}\) when
\(\sigma>1\), to \(X^{(1-\alpha)\omega_X}\) when \(\sigma=1\), and to
\(X^{1-\alpha}\) when \(\sigma<1\). Every exponent lies between zero and
one, so \(R'(X)\to\infty\). As \(X\to\infty\), marginal revenue converges to
zero from above if \(\sigma\geq1\); if \(\sigma<1\), it eventually
becomes negative. Hence \(R'(X)=1/B\) has exactly one solution. Operating profit
increases before this solution and decreases after it. The linear cost \(X/B\)
dominates revenue as \(X\to\infty\), so this solution is the unique global
maximum. Equation \eqref{eq:rewrite-monopoly-proof-curvature} also proves the
strict second-order condition:
\begin{equation}
 e_X(1-e_X)
 +\left(\alpha-\frac1\sigma\right)\left(1-\frac1\sigma\right)s_X(1-s_X)>0.
 \label{eq:rewrite-monopoly-soc}
\end{equation}
At the unique solution of $R'(X)=1/B$, $R''(X)<0$.
The implicit-function theorem therefore makes the maximizing service choice
smooth in $K,A,L,B$.
\end{proof}

\begin{lemma}[Optimal research expenditure]
\label{lem:rewrite-developer-verification}
Consider a feasible research trajectory $\{B_t,M_t\}_{t\geq0}$
for problem~\eqref{eq:rewrite-developer-reduced}, starting from
$B_0$ and satisfying the RSI technology
\eqref{eq:rewrite-capability-law}. Suppose $M_t>0$ and the
present value of net profit is finite. Suppose there exists
a continuously differentiable shadow-price trajectory
$\{q_t\}_{t\geq0}$ satisfying the research condition
\eqref{eq:rewrite-research-foc}, the costate equation
\eqref{eq:rewrite-costate}, and the transversality condition
\eqref{eq:rewrite-developer-tvc}.

Suppose $0<B_0<\overline B<\infty$ and, for every $t\geq0$
and every $b\in[B_0,\overline B)$,
\begin{equation}
 \begin{aligned}
 \pi_t(b)-\pi_t(B_t)
 &+\frac{1-\eta}{\eta}M_t
 \left[
 \left(\frac{b}{B_t}\right)^{\eta/(1-\eta)}-1
 \right]\\
 &\leq
 \frac{\overline B-B_t}{B_t}(U_t+M_t)
 \log\!\left(
 \frac{\overline B-B_t}{\overline B-b}
 \right),
 \end{aligned}
 \label{eq:rewrite-hamiltonian-support}
\end{equation}
where $\pi_t(b)$ is operating profit as defined in
equation~\eqref{eq:rewrite-service-choice}.
Then $\{M_t\}_{t\geq0}$ solves the developer's
problem~\eqref{eq:rewrite-developer-reduced}:
no other feasible research trajectory starting from
the same $B_0$ yields a higher present value of net profit.
\end{lemma}

\begin{proof}[Proof of Lemma~\ref{lem:rewrite-developer-verification}]
For the proof, define the transformed AI-efficiency state, or research depth, by
\begin{equation}
 R_B(B)=-\overline B\log\left(1-\frac{B}{\overline B}\right),
 \qquad
 B(R_B)=\overline B\left(1-e^{-R_B/\overline B}\right).
 \label{eq:rewrite-research-depth}
\end{equation}
It maps $B\in[B_0,\overline B)$ into
$R_B\in[R_{B,0},\infty)$, where $R_{B,t}=R_B(B_t)$ and
$R_{B,0}=R_B(B_0)$. Since $R_B'(B)=1/\psi(B)$, this transformation removes
the factor $\psi(B)$ from the RSI technology~\eqref{eq:rewrite-capability-law}:
\begin{equation}
 \dot R_B=\chi[B(R_B)M]^\eta.
 \label{eq:rewrite-research-depth-law}
\end{equation}
The corresponding current-value shadow price is $q_{R,t}=q_t\psi(B_t)$,
which is positive by the research condition~\eqref{eq:rewrite-research-foc}.
At each date $t$, hold $q_{R,t}$ at its value along the candidate trajectory
and define the maximized Hamiltonian for an alternative state $\widetilde R_B$:
\[
 \widehat{\mathcal H}_t(\widetilde R_B)
 \equiv\max_{\widetilde M\geq0}
 \left\{\pi_t(B(\widetilde R_B))-\widetilde M
 +q_{R,t}\chi[B(\widetilde R_B)\widetilde M]^\eta\right\}.
\]
Because $0<\eta<1$, the expression being maximized is strictly concave in
research expenditure. Using the candidate's research
condition~\eqref{eq:rewrite-research-foc}, the maximizing expenditure at
$\widetilde R_B$ is
\[
 M_t\left[\frac{B(\widetilde R_B)}{B_t}\right]^{\eta/(1-\eta)},
\]
where $M_t$ is the candidate's research expenditure. Substitution gives
\begin{equation}
 \widehat{\mathcal H}_t(\widetilde R_B)
 =\pi_t(B(\widetilde R_B))+
 \frac{1-\eta}{\eta}M_t
 \left[\frac{B(\widetilde R_B)}{B_t}\right]^{\eta/(1-\eta)}.
 \label{eq:rewrite-maximized-depth-hamiltonian}
\end{equation}
Differentiating the maximized Hamiltonian with respect to its state argument
and using the envelope condition~\eqref{eq:rewrite-profit-envelope} gives
\begin{equation}
 \widehat{\mathcal H}'_t(R_{B,t})
 =\frac{\psi(B_t)}{B_t}(U_t+M_t).
 \label{eq:rewrite-depth-hamiltonian-slope}
\end{equation}
Because
\[
 R_B(b)-R_{B,t}
 =\overline B\log\!\left(\frac{\overline B-B_t}{\overline B-b}\right),
\]
condition~\eqref{eq:rewrite-hamiltonian-support} is equivalent to the
following inequality for every $\widetilde R_B\geq R_{B,0}$:
\[
 \widehat{\mathcal H}_t(\widetilde R_B)
 \leq\widehat{\mathcal H}_t(R_{B,t})
 +\widehat{\mathcal H}'_t(R_{B,t})
 (\widetilde R_B-R_{B,t}).
\]

Let $D_t=\exp(-\int_0^t r_s\dd s)$ be the discount factor and
$\mu_t=D_tq_{R,t}$ the present-value shadow price of research depth.
The research condition~\eqref{eq:rewrite-research-foc} and the RSI
technology~\eqref{eq:rewrite-capability-law} imply $q_t\eta\dot B_t=M_t$.
Differentiating $q_{R,t}=q_t\psi(B_t)$ and using the costate
equation~\eqref{eq:rewrite-costate} therefore gives
\[
 \dot q_{R,t}=r_tq_{R,t}-\frac{\psi(B_t)}{B_t}(U_t+M_t),
 \qquad
 \dot\mu_t=-D_t\widehat{\mathcal H}'_t(R_{B,t}).
\]

Consider any feasible alternative $\{\widetilde B_t,\widetilde M_t\}_{t\geq0}$
starting from the same $B_0$, and write $\widetilde R_{B,t}=R_B(\widetilde B_t)$.
Here feasible expenditure is nonnegative and integrable on every finite
interval, with a locally absolutely continuous state satisfying the RSI
technology almost everywhere.
Let
\[
 J_T=\int_0^T D_t[\pi_t(B_t)-M_t]\dd t
\]
denote discounted net profit through date $T$ under the candidate, and define
$\widetilde J_T$ analogously for the alternative.
At each date, the alternative's Hamiltonian cannot exceed
$\widehat{\mathcal H}_t(\widetilde R_{B,t})$, whereas the candidate attains
$\widehat{\mathcal H}_t(R_{B,t})$. Applying
condition~\eqref{eq:rewrite-hamiltonian-support} and the transformed costate
equation yields, for almost every $t$,
\[
 \begin{aligned}
 &D_t\left[
 \pi_t(\widetilde B_t)-\widetilde M_t
 -\pi_t(B_t)+M_t
 \right]\\
 &\quad\leq D_t\widehat{\mathcal H}'_t(R_{B,t})
 (\widetilde R_{B,t}-R_{B,t})
 -\mu_t(\dot{\widetilde R}_{B,t}-\dot R_{B,t})\\
 &\quad=-\frac{\mathrm d}{\mathrm dt}
 \left[\mu_t(\widetilde R_{B,t}-R_{B,t})\right].
 \end{aligned}
\]
Integrating from zero to $T$ and using the common initial
state gives
\[
 \widetilde J_T-J_T
 \leq-\mu_T(\widetilde R_{B,T}-R_{B,T})
 \leq\mu_TR_{B,T},
\]
where the second inequality uses $\mu_T>0$ and $\widetilde R_{B,T}\geq0$.
Since
\[
 \psi(B_t)R_{B,t}
 =-(\overline B-B_t)\log\left(1-\frac{B_t}{\overline B}\right)
 \leq B_t,
\]
the developer's transversality condition~\eqref{eq:rewrite-developer-tvc} gives
\[
 0\leq\mu_TR_{B,T}
 =D_Tq_T\psi(B_T)R_{B,T}
 \leq D_Tq_TB_T\longrightarrow0.
\]
The candidate's present value is finite by assumption. Hence
\[
 \limsup_{T\to\infty}\widetilde J_T
 \leq\lim_{T\to\infty}J_T<\infty,
\]
so no feasible alternative yields a higher present value of net profit.
The argument requires a global upper tangent at the candidate state at each
date, not concavity at every counterfactual state.
\end{proof}

\Needspace{8\baselineskip}
\begin{lemma}[Sufficient conditions for Hamiltonian concavity]
\label{lem:rewrite-hamiltonian-concavity}
Consider the maximized Hamiltonian in
equation~\eqref{eq:rewrite-maximized-depth-hamiltonian} for a feasible
trajectory satisfying the research condition~\eqref{eq:rewrite-research-foc},
with $M_t>0$ and $0<B_0<\overline B<\infty$.
Let $X_t^*(b)$ denote the optimal AI-service choice in
problem~\eqref{eq:rewrite-service-choice} at efficiency $b$, and let
$\varepsilon_{B,t}(b)=\partial\ln X_t^*(b)/\partial\ln b$
denote its elasticity with respect to AI efficiency.
Suppose that, for every $t\geq0$ and $b\in[B_0,\overline B)$,
\begin{equation}
 \psi(b)[\varepsilon_{B,t}(b)-2]\leq\frac{b}{\overline B},
 \label{eq:rewrite-capped-concavity-gate}
\end{equation}
and
\begin{equation}
 \frac{B_0}{\overline B}\geq
 \max\left\{0,\frac{2\eta-1}{\eta}\right\}.
 \label{eq:rewrite-depth-research-concavity}
\end{equation}
Then the maximized Hamiltonian is concave in the transformed state
$R_B\in[R_B(B_0),\infty)$ at every date. Consequently, the support
condition~\eqref{eq:rewrite-hamiltonian-support} in
Lemma~\ref{lem:rewrite-developer-verification} holds.
\end{lemma}

\begin{proof}[Proof of Lemma~\ref{lem:rewrite-hamiltonian-concavity}]
Fix a date $t$. Lemma~\ref{lem:rewrite-monopoly-choice} ensures that
$X_t^*(b)$ is unique and smooth. Let $s_{X,t}(b)$ and $e_{X,t}(b)$ be the
corresponding elasticities defined in equations~\eqref{eq:rewrite-sx}
and~\eqref{eq:rewrite-inverse-demand-elasticity}. Implicit differentiation
of the monopoly condition~\eqref{eq:rewrite-monopoly-foc} gives
\begin{equation}
 \varepsilon_{B,t}(b)\equiv
 \frac{\partial\ln X_t^*(b)}{\partial\ln b}
 =\frac{1-e_{X,t}(b)}{e_{X,t}(b)[1-e_{X,t}(b)]
 +(\alpha-\sigma^{-1})(1-\sigma^{-1})
 s_{X,t}(b)[1-s_{X,t}(b)]}.
 \label{eq:rewrite-service-capability-elasticity}
\end{equation}
The denominator is strictly positive by
equation~\eqref{eq:rewrite-monopoly-soc}. The envelope
condition~\eqref{eq:rewrite-profit-envelope} gives
$\pi_t'(b)=X_t^*(b)/b^2$, so
\begin{equation}
 \pi_t''(b)=\frac{X_t^*(b)}{b^3}[\varepsilon_{B,t}(b)-2].
 \label{eq:rewrite-profit-curvature}
\end{equation}

The transformation~\eqref{eq:rewrite-research-depth} satisfies
$B'(R_B)=\psi(b)$ and $B''(R_B)=-\psi(b)/\overline B$, where $b=B(R_B)$.
The chain rule therefore gives
\begin{equation}
 \frac{\mathrm d^2\pi_t(B(R_B))}{\mathrm dR_B^2}
 =\frac{X_t^*(b)\psi(b)}{b^3}
 \left[\psi(b)(\varepsilon_{B,t}(b)-2)
 -\frac{b}{\overline B}\right].
 \label{eq:rewrite-capped-profit-curvature}
\end{equation}
Condition~\eqref{eq:rewrite-capped-concavity-gate} makes this derivative
nonpositive throughout the reachable state domain.

The research term in \eqref{eq:rewrite-maximized-depth-hamiltonian} is a
positive multiple of $B(R_B)^{\eta/(1-\eta)}$. Its curvature is
\begin{equation}
 \frac{\mathrm d^2 B(R_B)^{\eta/(1-\eta)}}{\mathrm dR_B^2}
 =\frac{\eta}{1-\eta}b^{\eta/(1-\eta)-2}\psi(b)
 \left[\frac{2\eta-1}{1-\eta}\psi(b)-\frac b{\overline B}\right].
 \label{eq:rewrite-depth-research-curvature}
\end{equation}
The bracket is nonpositive exactly when $b/\overline B\geq(2\eta-1)/\eta$.
Since $b\geq B_0$, condition~\eqref{eq:rewrite-depth-research-concavity}
makes the research term concave as well. The maximized Hamiltonian is
therefore concave, so its tangent at the candidate state is a global upper
bound. As shown in the proof of
Lemma~\ref{lem:rewrite-developer-verification}, this is
condition~\eqref{eq:rewrite-hamiltonian-support}.

If $\eta\leq1/2$, condition~\eqref{eq:rewrite-depth-research-concavity}
holds automatically. But for every fixed
$0<\eta<1$ its right-hand side is strictly below one, so it also holds
without this restriction when initial efficiency is sufficiently close to
the upper bound. These curvature tests are sufficient, not necessary: they
can fail even when the Hamiltonian has the required supporting tangent at
the candidate state.
\end{proof}

\Needspace{4\baselineskip}
\phantomsection
\label{app:rewrite-finite-existence}
\label{proof:rewrite-equilibrium-regimes}
\begin{proof}[Proof of Proposition~\ref{prop:rewrite-equilibrium-regimes}]

The proof has three tasks. First, Lemma~\ref{lem:rewrite-monopoly-choice}
gives a unique, smooth, globally optimal within-date service choice at given
states. I combine this choice with the other equilibrium equations to
identify candidate long-run values and verify their admissibility, including
positive consumption. Finding these values does not yet establish the
existence of a trajectory approaching them.

Second, I write the exact dynamic equations in coordinates with finite
candidate limits. The linearization identifies the stable directions; I
also check that they can accommodate every sufficiently nearby initial
state. Counting stable directions alone would not establish this property.
The stable-manifold and inverse-function theorems then select the initial
consumption and shadow price for those states. They establish convergent
solutions of the nonlinear equations for all future time, not merely
solutions of their linear approximation.

Third, I reconstruct the original variables and verify feasibility,
market clearing, and both transversality conditions. Household concavity
and Lemma~\ref{lem:rewrite-developer-verification} establish optimality.
For the developer, I verify the curvature conditions of
Lemma~\ref{lem:rewrite-hamiltonian-concavity} at every AI-efficiency level
reachable under alternative research policies from the initial efficiency.
The upper bound makes this possible for initial efficiency sufficiently
close to it. These final checks turn the constructed solutions of the
first-order conditions into equilibria. The result is local in initial
states, not in the time horizon.

\par\medskip\noindent\textit{Step 1. Feasible efficiency and candidate long-run values.}

Let \(D=\overline B-B\). The state equation with an upper bound gives
\[
 \frac{\dot D}{D}=-\frac{\chi}{\overline B}(BM)^\eta.
\]
Integration gives
\begin{equation}
 \overline B-B_t=(\overline B-B_0)
 \exp\left\{-\frac{\chi}{\overline B}
 \int_0^t(B_sM_s)^\eta\dd s\right\}>0.
 \label{eq:rewrite-cap-gap}
\end{equation}
Local integrability makes the exponent finite at every finite date, so
\(D_t>0\).

For the candidate limits, impose $B=\overline B$ only in the normalized
boundary equations. Actual trajectories will satisfy $B<\overline B$
at every finite date.


Write $\theta=(1-\alpha)/\alpha$. For $\sigma\ne1$, let
$\varphi=(\sigma-1)/\sigma$ and define
\[
 m(s)=1-\frac{1-s}{\sigma}-\alpha s,\qquad
 \mathcal H_\sigma(s)=
 \left[(1-\alpha)s\,m(s)
       \left(\frac{\omega_X}{s}\right)^{1/\varphi}\right]^\theta.
\]
Here $m(s)=1-e_X$ is the marginal-revenue factor, not research compute.
The static monopoly optimum has $m(s_X)>0$ by
Lemma~\ref{lem:rewrite-monopoly-choice}. The CES share, monopoly condition,
and production function, equations~\eqref{eq:rewrite-eq-sx},
\eqref{eq:rewrite-eq-monopoly}, and \eqref{eq:rewrite-eq-y}, imply
\begin{align}
 \frac{X}{AN}
 &=\left[\frac{\omega_Ls_X}{\omega_X(1-s_X)}\right]^{1/\varphi},
 &
 \frac{U}{Y}&=(1-\alpha)s_Xm(s_X), \label{eq:rewrite-frontier-static-shares}\\
 \frac{Y}{K}&=B^\theta\mathcal H_\sigma(s_X).
 \label{eq:rewrite-frontier-share-map}
\end{align}
The identities $Z=X(\omega_X/s_X)^{1/\varphi}$ and
$X=B(1-\alpha)s_Xm(s_X)Y$ give the output--capital ratio in
\eqref{eq:rewrite-frontier-share-map}.

If $0<\sigma<1$, positive marginal revenue requires
\[
 s>s_{\min}\equiv\frac{1-\sigma}{1-\alpha\sigma}.
\]
On $(s_{\min},1)$, both $m(s)$ and $s^{-1/(\sigma-1)}$ increase strictly.
Since
\[
 \mathcal H_\sigma(s)^{1/\theta}
 =(1-\alpha)\omega_X^{1/\varphi}
 m(s)s^{-1/(\sigma-1)},
\]
the share map increases from zero to
$\mathcal H_\sigma(1)=[(1-\alpha)^2\omega_X^{1/\varphi}]^\theta$.
Thus the equation
\begin{equation}
 \frac{\rho+\gamma+\delta}{\alpha}
 =\overline B^\theta\mathcal H_\sigma(\overline s_X)
 \label{eq:rewrite-terminal-share-root}
\end{equation}
has exactly one interior solution if and only if
$\overline B>\mathcal B(\sigma)$.

If $\sigma>1$, the share map decreases strictly on $(0,1)$.
To verify the sign, put $a=\sigma^{-1}-\alpha$, so $m(s)=\varphi+as$.
Its logarithmic derivative has the sign of $a(\sigma-2)s-\varphi$.
If $a(\sigma-2)\leq0$ this is negative. If $\sigma>2$ and $a>0$,
$a(\sigma-2)<(\sigma-2)/\sigma<\varphi$; if $1<\sigma<2$ and $a<0$,
$a(\sigma-2)<(1-\sigma^{-1})(2-\sigma)<\varphi$.
The map therefore decreases from infinity to $\mathcal H_\sigma(1)$.
Equation~\eqref{eq:rewrite-terminal-share-root} has exactly one interior
solution if and only if $\overline B<\mathcal B(\sigma)$.
Since $\mathcal R(B)=\alpha B^\theta\mathcal H_\sigma(1)-\delta$,
these calculations also derive
\eqref{eq:rewrite-ai-terminal-ratio}--\eqref{eq:rewrite-critical-frontier}.
The interior-share condition is therefore
$\mathcal R(\overline B)>\rho+\gamma$ for $\sigma<1$, but
$\mathcal R(\overline B)<\rho+\gamma$ for $\sigma>1$.

\Needspace{4\baselineskip}
For either interior-share case, let $x,y,u,k,c$ denote $X,Y,U,K,C$ divided
by $AN$, respectively. Construct the normalized terminal values sequentially:
\begin{align}
 \overline x&=
 \left[\frac{\omega_L\overline s_X}
 {\omega_X(1-\overline s_X)}\right]^{1/\varphi},
 &\overline u&=\frac{\overline x}{\overline B},\nonumber\\
 \overline y&=
 \frac{\overline u}{(1-\alpha)\overline s_Xm(\overline s_X)},
 &\overline k&=\frac{\alpha\overline y}{\rho+\gamma+\delta},\nonumber\\
 \overline c&=\overline y-\overline u
       -(\delta+n+\gamma)\overline k,
 &\overline r&=\rho+\gamma.
 \label{eq:rewrite-labor-terminal-values}
\end{align}
At $\sigma=1$, define $\beta=(1-\alpha)\omega_X$.
The exact static solution is
\[
 y=\left[(\beta^2B)^\beta k^\alpha\right]^{1/(1-\beta)},
 \qquad u=\beta^2y,\qquad x=Bu.
\]
The equation $y/k=(\rho+\gamma+\delta)/\alpha$ has exactly one positive
solution for every finite $\overline B>0$:
\begin{equation}
 \overline k=
 \left[(\beta^2\overline B)^\beta
 \left(\frac{\rho+\gamma+\delta}{\alpha}\right)^{\beta-1}
 \right]^{1/(1-\alpha-\beta)}.
 \label{eq:rewrite-unit-terminal-capital}
\end{equation}
Set
$\overline y=[(\rho+\gamma+\delta)/\alpha]\overline k$,
$\overline u=\beta^2\overline y$,
and $\overline x=\overline B\,\overline u$, with $\overline c,\overline r$
as in \eqref{eq:rewrite-labor-terminal-values}.

Consumption positivity does not require an additional restriction.
Positive marginal revenue implies $0<m(\overline s_X)<1$.
Since $0<\overline s_X<1$, in all three cases
\[
 \frac{\overline u}{\overline y}
 =(1-\alpha)\overline s_Xm(\overline s_X)<1-\alpha,\qquad
 \frac{(\delta+n+\gamma)\overline k}{\overline y}
 =\frac{\alpha(\delta+n+\gamma)}{\rho+\gamma+\delta}<\alpha.
\]
It follows that $\overline c/\overline y>0$.

\par\medskip\noindent\textit{Step 2. The three positive-labor-share cases.}

The cases in Proposition~\ref{prop:rewrite-capped-complements-existence},
\ref{prop:rewrite-capped-unit-existence}, and
\ref{prop:rewrite-capped-substitutes-labor-existence} share one
normalization. Divide capital, consumption, and the AI-efficiency shadow
price by effective labor, and multiply the remaining efficiency gap by
effective labor:
\begin{equation}
 k=\frac K{AN},\quad c=\frac C{AN},\quad
 d=(\overline B-B)AN,\quad \widetilde q=\frac q{AN},
 \quad \tau=(AN)^{-1}.
 \label{eq:rewrite-labor-proof-coordinates}
\end{equation}
Dividing by $AN$ removes effective-labor growth. Multiplying the shrinking
gap by $AN$ allows it to converge to a positive value:
$\dot d/d=n+\gamma-\dot B/(\overline B-B)$.
The variable $d$ measures the scaled gap, not its rate of closure.
The constants $\overline k,\overline d$ are the limiting values of $k,d$;
they are determined by the model, not additional parameters.
Also $\dot\tau=-(n+\gamma)\tau$, so $\tau\to0$.
These coordinates leave the model unchanged: $K=k/\tau$ and
$B=\overline B-d\tau$ at every finite date.

Given $k,B$, solve the unique static monopoly
problem for $y,x,u$, and set $r=\alpha y/k-\delta$.
Define $a=\dot B/(\overline B-B)$, the proportional rate of closure of the
AI-efficiency gap. The research condition gives
\[
 B=\overline B-d\tau,\qquad
 M=\left[\widetilde qd\,\chi\eta\frac{B^\eta}{\overline B}
       \right]^{1/(1-\eta)},\qquad
 a=\chi\frac{B^\eta M^\eta}{\overline B}.
\]
Substituting \eqref{eq:rewrite-eq-research} and
\eqref{eq:rewrite-equilibrium-dynamics} gives the exact dynamic equations
\begin{subequations}
\label{eq:rewrite-labor-proof-dynamics}
\begin{align}
 \frac{\dot k}{k}&=\frac{y-c-u-M\tau}{k}-\delta-n-\gamma,\\
 \frac{\dot c}{c}&=r-\rho-\gamma,\\
 \frac{\dot d}{d}&=n+\gamma-a,\\
 \frac{\dot{\widetilde q}}{\widetilde q}
 &=r-\frac{x}{\widetilde qB^2}
   -\eta\frac{ad\tau}{B}+a-n-\gamma,\\
 \dot\tau&=-(n+\gamma)\tau.
\end{align}
\end{subequations}
In particular, population growth cancels from the normalized consumption
equation; it is not subtracted a second time.

Complete the terminal construction by setting
\begin{equation}
 \overline M=
 \left[\frac{(n+\gamma)\overline B^{1-\eta}}{\chi}\right]^{1/\eta},
 \qquad
 \overline{\widetilde q}=\frac{\overline x}{\overline r\,\overline B^2},
 \qquad
 \overline d=\frac{\overline M}{\eta(n+\gamma)\overline{\widetilde q}}.
 \label{eq:rewrite-labor-terminal-research}
\end{equation}
All are positive. Direct substitution shows that
$(\overline k,\overline c,\overline d,\overline{\widetilde q},0)$
is a fixed point of \eqref{eq:rewrite-labor-proof-dynamics}.
This boundary point is a device for constructing paths with $\tau>0$ and
$B<\overline B$, not an admissible date-zero state itself.

I next verify saddle-path stability. Write $\mu=\eta/(1-\eta)$,
$r_k=\partial r/\partial\log k$, and
$a_k=\partial[(y-u)/k]/\partial\log k+\overline c/\overline k$,
evaluated at the fixed point. In
$(\log k,\log c,\log d,\log\widetilde q,\tau)$ coordinates the Jacobian has
two diagonal blocks: one for capital and consumption, and one for the scaled
efficiency gap and normalized shadow price,
\begin{equation}
 J_R=\begin{pmatrix}
 a_k&-\overline c/\overline k\\ r_k&0
 \end{pmatrix},
 \qquad
 J_B=\begin{pmatrix}
 -(n+\gamma)\mu&-(n+\gamma)\mu\\
 (n+\gamma)\mu&\overline r+(n+\gamma)\mu
 \end{pmatrix},
 \label{eq:rewrite-labor-proof-blocks}
\end{equation}
and the additional eigenvalue $-(n+\gamma)$. There can be forcing from
$\tau$ to both blocks and from $(k,c)$ to $(d,\widetilde q)$, but not in the
reverse directions at the fixed point.

To check $r_k<0$, let $s=s_X$, $u_\sigma=1/\sigma$, $e=e_X$, and let $D(s)$
be the positive curvature expression in
\eqref{eq:rewrite-monopoly-proof-curvature}. Implicit differentiation of the
monopoly condition gives
\[
 \frac{\partial\log X}{\partial\log K}
 =\frac{\alpha(1-e)}{D(s)},\qquad
 r_k=(r+\delta)(1-\alpha)
 \left[\frac{\alpha s(1-e)}{D(s)}-1\right].
\]
Moreover,
\[
 D(s)-\alpha s(1-e)
 =(1-s)\{u_\sigma(1-e)+
       (\alpha-u_\sigma)(1-u_\sigma)s\}>0.
\]
The sign is immediate when $u_\sigma\geq1$ or $u_\sigma\leq\alpha$.
For $\alpha<u_\sigma<1$, the braces are linear in $s$, with positive
endpoint values $u_\sigma(1-u_\sigma)$ and
$u_\sigma^2+\alpha(1-2u_\sigma)$.
The latter is positive if $u_\sigma\leq1/2$; otherwise $\alpha<u_\sigma$
implies it exceeds $u_\sigma(1-u_\sigma)>0$.
Thus $r_k<0$ in every interior-share case.

Both blocks have negative determinants:
\[
 \det J_R=(\overline c/\overline k)r_k<0,\qquad
 \det J_B=-(n+\gamma)\mu\overline r<0.
\]
There are therefore three stable and two unstable eigenvalues. The stable
direction in $J_R$ has a nonzero $k$ component, and the stable direction in
$J_B$ has a nonzero $d$ component. To verify that their number suffices to
select the jump variables, consider a vector in the stable subspace whose
state components $(k,d,\tau)$ are all zero. Its linearized $\tau$ path is
zero. The unforced capital--consumption block must then follow its stable direction;
its zero initial $k$ component forces its $c$ component to be zero as well.
With that forcing absent, the same argument in the research block forces
$\widetilde q=0$. The projection of the three-dimensional stable subspace
onto the three predetermined coordinates is therefore injective and hence
invertible. This argument also covers repeated stable eigenvalues.

Lemma~\ref{lem:rewrite-monopoly-choice} makes the static solution smooth.
The normalized vector field extends smoothly through $\tau=0$, since
$B=\overline B-d\tau$ remains positive nearby. The stable-manifold theorem
and the inverse-function theorem therefore give a local graph
$(c,\widetilde q)$ over $(k,d,\tau)$; see \citet{dyatlov2018}.
Choose a small neighborhood of the full fixed point and then a smaller
initial-state neighborhood $\mathcal U_L$ of
$(\overline k,\overline d,0)$. The size of this neighborhood depends on the parameters
and the case under consideration. In original initial stocks the restriction is
\begin{equation}
 (k_0,d_0,\tau_0)=
 \left(\frac{K_0}{A_0N_0},\,
       (\overline B-B_0)A_0N_0,\,
       \frac{1}{A_0N_0}\right)\in\mathcal U_L.
 \label{eq:rewrite-labor-local-initial}
\end{equation}
Closeness is required in all three coordinates; $B_0$ close to
$\overline B$ alone is not sufficient. The point with $\tau=0$
is a long-run limit, not a finite-date initial state.
Every physical initial state satisfying
\eqref{eq:rewrite-labor-local-initial} has a solution on this
graph that stays in the full neighborhood for all future time and converges
to the fixed point. It is unique among solutions with those properties.
Positivity of $\tau,d,c,\widetilde q,k$ and the original RSI technology
ensure $0<B_t<\overline B$ and $M_t>0$ at every finite date.
Reconstruction gives the original initial stocks, the equilibrium equations,
and the limiting rates in
\eqref{eq:rewrite-complements-terminal}--\eqref{eq:rewrite-substitutes-labor-terminal}.
Indeed, the normalized vector field tends to zero at its fixed point, so
the time derivatives of its smooth static ratios tend to zero as well.
This step, not convergence of ratios alone, establishes convergence of
the growth rates.
Optimality and transversality are verified below.

\par\medskip\noindent\textit{Step 3. The AI-dominated case.}

Let $\sigma>1$ and $\overline B>\mathcal B(\sigma)$.
Here $\varphi=(\sigma-1)/\sigma\in(0,1)$.
I seek an equilibrium in which capital grows faster than effective labor.
Along such a trajectory, $K/(AN)$ would diverge, so I normalize consumption
and the shadow value by capital and scale the remaining AI-efficiency gap
by capital:
\[
 h=(AN/K)^\varphi,\quad d=(\overline B-B)K,\quad
 c=C/K,\quad v=q/K,\quad \tau=(AN)^{-1}.
\]
Within this construction, $y,x,u$ denote $Y/K,X/K,U/K$, respectively.
The constant $\overline d$ now denotes the limit of $(\overline B-B)K$,
not of $(\overline B-B)AN$; no original state or control is redefined.
Then
\[
 B=\overline B-d\tau h^{1/\varphi},\quad
 M=\left[vd\,\chi\eta B^\eta/\overline B\right]^{1/(1-\eta)},\quad
 a=\chi B^\eta M^\eta/\overline B,\quad r=\alpha y-\delta.
\]
Capital's growth rate is $g_K=y-c-u-M\tau h^{1/\varphi}-\delta$.
The exact equations are
\begin{subequations}
\label{eq:rewrite-ai-proof-dynamics}
\begin{align}
 \dot h&=\varphi h(n+\gamma-g_K),\\
 \dot d/d&=g_K-a,\\
 \dot c/c&=n+r-\rho-g_K,\\
 \dot v/v&=r-x/(vB^2)-\eta ad\tau h^{1/\varphi}/B+a-g_K,\\
 \dot\tau&=-(n+\gamma)\tau.
\end{align}
\end{subequations}
The CES term divided by capital is
$[\omega_Lh+\omega_Xx^\varphi]^{1/\varphi}$.
At $h=0$ the monopoly equation has the strictly positive curvature
$D(1)=\alpha(1-\alpha)$. Its implicit static solution is therefore smooth
in $(h,B)$ near $(0,\overline B)$. For the remaining terms, replacing
$h^{1/\varphi}$ by $\max\{h,0\}^{1/\varphi}$ outside the physical domain
gives a $C^1$ extension because $1/\varphi>1$. The $C^1$ stable-manifold
theorem applies; see \citet[p.~7]{feldman2021stable}. The extension is only a
mathematical device and does not change any equation on $h\geq0$.

The limiting interest and aggregate growth rates are
\begin{equation}
 \overline r=\mathcal R(\overline B),\qquad
 \overline g=n+\overline r-\rho>n+\gamma.
 \label{eq:rewrite-substitutes-terminal-growth}
\end{equation}
At $h=\tau=0$, the terminal constants are
\begin{align}
 \overline y&=\frac{\mathcal R(\overline B)+\delta}{\alpha},&
 \overline u&=(1-\alpha)^2\overline y,&
 \overline x&=\overline B\,\overline u,\nonumber\\
 \overline r&=\mathcal R(\overline B),&
 \overline g&=n+\overline r-\rho,&
 \overline c&=\alpha(1-\alpha)\overline y+\rho-n>0,\nonumber\\
 \overline v&=\overline x/(\overline r\,\overline B^2),&
 \overline M&=
 \left[\overline g\,\overline B^{1-\eta}/\chi\right]^{1/\eta},&
 \overline d&=\overline M/(\eta\overline g\,\overline v).
 \label{eq:rewrite-ai-terminal-values}
\end{align}
The threshold implies $\overline g>n+\gamma>0$ and
$\overline r>\rho+\gamma>0$, so all constants are positive.
Substitution verifies the fixed point
$(0,\overline d,\overline c,\overline v,0)$.

The $h$ and $\tau$ eigenvalues are
$\varphi(n+\gamma-\overline g)<0$ and $-(n+\gamma)<0$.
The first equals $\varphi[\rho+\gamma-\mathcal R(\overline B)]$.
It becomes zero at $\mathcal R(\overline B)=\rho+\gamma$, which is why
this stable-manifold argument does not cover the threshold.
The consumption-ratio eigenvalue is $\overline c>0$.
The remaining block in $(\log d,\log v)$ is
\[
 \begin{pmatrix}
 -\overline g\mu&-\overline g\mu\\
 \overline g\mu&\overline r+\overline g\mu
 \end{pmatrix},
 \qquad \det=-\overline g\mu\overline r<0.
\]
The linearized system is triangular in the order
$(h,\tau),c,(d,v)$. A stable vector with zero state components $(h,d,\tau)$
has zero $h,\tau$ paths; its consumption component must be zero because the
unforced consumption equation is unstable. The research block then has zero
$d$ component and hence zero $v$ component, as in the previous construction.
Thus the three-dimensional stable subspace projects invertibly onto
$(h,d,\tau)$. The same graph argument gives a neighborhood
$\mathcal U_A$ of $(0,\overline d,0)$ such that the initial-state condition is
\begin{equation}
 \left[\left(\frac{A_0N_0}{K_0}\right)^\varphi,\;
       (\overline B-B_0)K_0,\; \frac{1}{A_0N_0}\right]\in\mathcal U_A.
 \label{eq:rewrite-ai-local-initial}
\end{equation}
It constructs a convergent trajectory for every admissible initial state
in this neighborhood, with a unique nearby convergent choice
of $C_0,q_0$ and positive quantities at every finite date.

The normalized vector field tends to zero, so the time derivatives of the
smooth positive ratios $Y/K$, $X/K$, and $C/K$ tend to zero.
Consequently $g_Y,g_X,g_C,g_K\to\overline g$ and $s_X\to1$.
Log differentiation of the CES share in
equation~\eqref{eq:rewrite-eq-sx} and the wage equation
\eqref{eq:rewrite-eq-w} gives the exact identity
\[
 g_w=g_Y-n+\frac{\sigma-1}{\sigma}s_X(n+\gamma-g_X).
\]
Taking limits yields
\[
 g_w\to\overline g-n+\frac{\sigma-1}{\sigma}(n+\gamma-\overline g)
 =\gamma+\frac{\overline r-\rho-\gamma}{\sigma}.
\]
This proves the growth and distribution limits in
Proposition~\ref{prop:rewrite-capped-substitutes-existence}.

\par\medskip\noindent\textit{Step 4. Optimality and transversality.}

In either construction let $\overline g$ denote the terminal aggregate
growth rate: $n+\gamma$ in the positive labor-share case and
$n+\overline r-\rho$ in the AI-dominated case. The constructions imply
\[
 a\to\overline g,\qquad
 M\to\left[\frac{\overline g\,\overline B^{1-\eta}}{\chi}\right]^{1/\eta}>0,
 \qquad g_q\to\overline g,\qquad g_B\to0,\qquad M/Y\to0.
\]
Both transversality expressions in Definition~\ref{def:rewrite-equilibrium}
have limiting logarithmic growth
\begin{equation}
 -\rho+n+g_K-g_C\to n-\rho<0,\qquad
 -r+g_q+g_B\to\overline g-\overline r=n-\rho<0.
 \label{eq:rewrite-finite-existence-tvcs}
\end{equation}
They therefore converge to zero, rather than merely remaining bounded.

Household utility is finite because $\log(C/N)$ grows at most linearly and
$\rho>n$, as imposed in the household problem. Log utility is concave, and the
household budget is linear at given prices and distributed profits. More explicitly,
let $\lambda_t=e^{-\rho t}N_t/C_t$, so $\dot\lambda_t=-r_t\lambda_t$ by
\eqref{eq:rewrite-euler}. Concavity of log utility and equal initial capital
give, for any admissible alternative $(\widetilde C,\widetilde K)$,
\[
 \int_0^T e^{-\rho t}N_t
 \log\frac{\widetilde C_t}{C_t}\,\dd t
 \leq\int_0^T\lambda_t(\widetilde C_t-C_t)\,\dd t
 =-\lambda_T(\widetilde K_T-K_T)\leq\lambda_TK_T\to0.
\]
Here $\widetilde K_T\geq0$ and the last limit is the household TVC.
Final-good optimality follows from concavity and constant returns.
The pointwise monopoly optimum is global by
Lemma~\ref{lem:rewrite-monopoly-choice}.

For the dynamic developer, feasibility and first-order conditions still
require a sufficiency argument. Let $S_t= A_tN_t$ in the positive-share
construction and $S_t=K_t$ in the AI-dominated construction, and let
$D_t=\exp(-\int_0^t r_s\,\dd s)$. Homogeneity of the static problem and
compactness of the normalized state neighborhood give a constant $C_\pi$
such that $0\leq\pi_t(b)\leq C_\pi S_t$ for every date and every
$b\in[B_0,\overline B]$. Since
$\mathrm d\log(D_tS_t)/\mathrm dt\to\overline g-\overline r=n-\rho<0$,
$\int_0^\infty D_tS_t\,\dd t<\infty$. Thus discounted operating profit
is uniformly bounded above across all alternative research policies.
The candidate's present value of net profit is finite, because its research compute is
bounded and $r_t\to\overline r>0$. An alternative with infinite discounted
research expenditure has payoff $-\infty$ and cannot improve on it.

It remains to verify the curvature conditions of
Lemma~\ref{lem:rewrite-hamiltonian-concavity},
\eqref{eq:rewrite-capped-concavity-gate} and
\eqref{eq:rewrite-depth-research-concavity} for every date and every alternative
AI-efficiency level in $[B_0,\overline B)$, not just on the selected path.
First fix a compact normalized-state neighborhood and an interval of
alternative AI-efficiency levels $[b_*,\overline B]$, with $0<b_*<\overline B$.
In the positive labor-share construction, $K/(AN)$ stays in a compact
neighborhood of $\overline k>0$. Smoothness of the static monopoly solution
bounds $\varepsilon_{B,t}(b)$ uniformly for these states and for AI-efficiency levels near
$\overline B$. In the AI-dominated construction, $h$ stays near zero, and
the same uniform bound follows from the smooth static limit
$\varepsilon_{B,t}(b)\to1/\alpha$. Thus in either case there is a finite constant
$E$ with $\varepsilon_{B,t}(b)-2\leq E$ throughout that state domain.
This bound is chosen before shrinking the initial-state neighborhood.
Because $\mu=\eta/(1-\eta)$ is finite for each $0<\eta<1$, choose
$b_{**}\in[b_*,\overline B)$ sufficiently close to $\overline B$ that
\[
 (1-b_{**}/\overline B)\max\{E,\mu-1,0\}<b_{**}/\overline B.
\]
Shrink the open neighborhood $\mathcal U_L$ or $\mathcal U_A$ so that
$B_0>b_{**}$.
This is possible because $B_0=\overline B-d_0\tau_0$ in the first
construction and $B_0=\overline B-d_0\tau_0h_0^{1/\varphi}$ in the second.
Every feasible research policy has $b\geq B_0$, so the two curvature
conditions hold throughout the entire alternative-state domain at every
date. Lemma~\ref{lem:rewrite-hamiltonian-concavity} therefore gives the
support inequality~\eqref{eq:rewrite-hamiltonian-support}.
Lemma~\ref{lem:rewrite-developer-verification} and the verified TVC
then establish global developer optimality without imposing
$\eta\leq1/2$. The curvature conditions follow from the primitives and the
chosen initial-state neighborhood.

All quantities remain finite at every finite date, and market clearing and
the original initial conditions hold by reconstruction. The paths are
therefore equilibria, not merely solutions of the
canonical equations. The result is local only in its set of initial
stocks. For $\sigma>1$, at $\overline B=\mathcal B(\sigma)$ the interior-share terminal
point reaches the boundary and a stable eigenvalue in the boundary
normalization becomes zero; this proof does not apply. For
$\sigma<1$ and $\overline B\leq\mathcal B(\sigma)$, it does not rule out
equilibria with different long-run ratios. Nor does it imply existence at
arbitrary initial stocks or in the uncapped limit.

\par\medskip\noindent\textit{Step 5. Initial-state sets and the four conclusions.}

The preimage of each open neighborhood under its continuous
initial-state map, restricted to $A_0,N_0,K_0>0$ and
$0<B_0<\overline B$, is open in the original initial stocks.
It is nonempty. In the positive-labor-share cases, choose
$k_0,d_0$ close to $\overline k,\overline d$ and $\tau_0>0$
sufficiently small. Then $A_0N_0=1/\tau_0$,
$K_0=k_0/\tau_0$, and $B_0=\overline B-d_0\tau_0$ are admissible.
In the AI-dominated case, choose $d_0$ near $\overline d$ and
$h_0,\tau_0>0$ sufficiently small. The inverse relations are
$A_0N_0=1/\tau_0$, $K_0=1/(\tau_0h_0^{1/\varphi})$, and
$B_0=\overline B-d_0\tau_0h_0^{1/\varphi}$.
Any positive factorization of $A_0N_0$ supplies $A_0,N_0$.
These choices also satisfy the strict efficiency restriction in Step 4.

For each initial state in the resulting set, the local graph selects
$C_0,q_0>0$. These jump values are unique among the trajectories
that remain near the indicated normalized limit and converge to it.
The argument does not assert global uniqueness.

For $\sigma<1$ with $\overline B>\mathcal B(\sigma)$, and for
$\sigma>1$ with $\overline B<\mathcal B(\sigma)$, Step 1 gives a unique
interior $\overline s_X$. Step 2 then yields
$g_Y-n\to\gamma$, $g_w\to\gamma$, and $r\to\rho+\gamma$.
At $\sigma=1$, the exact CES limit gives $s_X=\omega_X$ at every date,
and the same construction works for every finite positive upper bound.
These establish the first three cases of the proposition.
Step 3 gives $s_X\to1$, $r\to\mathcal R(\overline B)$,
$g_Y-n\to\mathcal R(\overline B)-\rho$, and the wage-growth limit in
\eqref{eq:rewrite-substitutes-wage-growth}, establishing the fourth case.
Step 4 verifies the agent-optimality and transversality requirements of
Definition~\ref{def:rewrite-equilibrium} for all these trajectories.
\end{proof}

%\Needspace{20\baselineskip}
%\begin{proposition}[Wage growth and labor's income share]
%\label{prop:rewrite-output-wage-wedge}
%Consider the long-run equilibria constructed in Proposition~\ref{prop:rewrite-equilibrium-regimes}. Then
%\[
% \lim_{t\to\infty}g_w>\gamma
% \quad\Longleftrightarrow\quad
% \lim_{t\to\infty}\frac{wL}{Y}=0,
%\]
%and either condition holds precisely when the equilibrium is AI-dominated. Along such an equilibrium, output per person grows asymptotically faster than the real wage:
%\begin{equation}
% \lim_{t\to\infty}\bigl[(g_Y-n)-g_w\bigr]
% =\left(1-\frac{1}{\sigma}\right)
%  (\overline r-\rho-\gamma)>0.
% \label{eq:rewrite-output-wage-wedge}
%\end{equation}
%Because $L=N$, the same wedge is the asymptotic rate at which labor's income share falls:
%\begin{equation}
% \lim_{t\to\infty}\frac{\mathrm d}{\mathrm dt}
% \log\left(\frac{wL}{Y}\right)
% =-\left(1-\frac{1}{\sigma}\right)
%  (\overline r-\rho-\gamma)<0.
% \label{eq:rewrite-labor-share-decline-rate}
%\end{equation}
%Equivalently,
%\begin{equation}
% -\lim_{t\to\infty}\frac{\mathrm d}{\mathrm dt}
%  \log\left(\frac{wL}{Y}\right)
% =(\sigma-1)\lim_{t\to\infty}(g_w-\gamma)>0.
% \label{eq:rewrite-wage-premium-share-decline}
%\end{equation}
%\end{proposition}

%\Needspace{4\baselineskip}
%\phantomsection
%\label{proof:rewrite-output-wage-wedge}
%\begin{proof}[Proof of Proposition~\ref{prop:rewrite-output-wage-wedge}]
%The first three cases of Proposition~\ref{prop:rewrite-equilibrium-regimes} give $g_w\to\gamma$ and a positive limiting labor share, using $wL/Y=(1-\alpha)(1-s_X)$. Proposition~\ref{prop:rewrite-capped-substitutes-existence} implies both $g_w\to\gamma+(\overline r-\rho-\gamma)/\sigma>\gamma$ and $wL/Y\to0$. This proves the equivalence. Subtract the wage-growth limit in \eqref{eq:rewrite-substitutes-wage-growth} from the per-person output-growth limit in \eqref{eq:rewrite-substitutes-ai-limits}. Moreover, $\mathrm d\log(wL/Y)/\mathrm dt=g_w+n-g_Y$ because $L=N$. Finally, combine \eqref{eq:rewrite-labor-share-decline-rate} with \eqref{eq:rewrite-substitutes-wage-growth}.
%\end{proof}

\Needspace{7\baselineskip}
\begin{proof}[Proof of Proposition~\ref{prop:rewrite-uncapped-complements-bounds}]
\phantomsection
\label{proof:rewrite-uncapped-complements-bounds}
Equations~\eqref{eq:rewrite-uncapped-complements-z-bound}
and~\eqref{eq:rewrite-uncapped-complements-y-bound} follow because
$\sigma/(\sigma-1)<0$. Define $k=K/(AL)$. The resource constraint
\eqref{eq:rewrite-eq-kdot}, nonnegative consumption and compute expenditure,
and $L=N$ imply
\begin{equation}
 \dot k\leq
 \omega_L^{\sigma(1-\alpha)/(\sigma-1)}k^\alpha
 -(\delta+n+\gamma)k.
 \label{eq:rewrite-uncapped-complements-k-bound}
\end{equation}
Since $\delta+n+\gamma>0$ and $0<\alpha<1$, the right-hand side is
negative above a finite positive level. In particular,
\[
 k_t\leq
 \max\left\{k_0,
 \left[\frac{\omega_L^{\sigma(1-\alpha)/(\sigma-1)}}
 {\delta+n+\gamma}\right]^{1/(1-\alpha)}\right\}.
\]
Equation~\eqref{eq:rewrite-uncapped-complements-y-bound} then bounds
$Y/(AL)$ uniformly. Since $A_t=A_0e^{\gamma t}$ and $L=N$,
$Y/N$ is bounded above by a constant times $e^{\gamma t}$.
Taking logarithms, dividing by $t$, and taking the upper limit gives
\eqref{eq:rewrite-uncapped-complements-growth-bound}.
\end{proof}

% Additional results retained for possible later use; not part of the proposition.
% The same bounds keep $K$ and $Y$ finite on every bounded time interval.
% For any such horizon $T$, integrating the resource constraint gives
% \[
%  \int_0^T M_t\dd t\leq K_0+\int_0^T Y_t\dd t<\infty.
% \]
% The uncapped RSI technology and H\"older's inequality imply
% \begin{align}
%  B_T^{1-\eta}
%  &=B_0^{1-\eta}+(1-\eta)\chi\int_0^T M_t^\eta\dd t
%  \nonumber\\
%  &\leq B_0^{1-\eta}
%  +(1-\eta)\chi T^{1-\eta}
%  \left(\int_0^T M_t\dd t\right)^\eta<\infty.
%  \label{eq:rewrite-uncapped-complements-b-bound}
% \end{align}
% Thus $B$ cannot diverge at a finite date either. This argument bounds
% stocks and output; it is not a regularity theorem for all instantaneous
% control choices or an equilibrium-existence proof.
%
% There is also no finite-horizon unbounded-payoff argument of the kind used
% for substitution in Proposition~\ref{prop:rewrite-research-scale}.
% At given positive continuous paths of $K,A,L$ on $[0,T]$, optimized operating
% profit satisfies
% \[
%  0\leq\pi_t(B)\leq p_XX\leq(1-\alpha)Y
%  \leq(1-\alpha)\omega_L^{\sigma(1-\alpha)/(\sigma-1)}
%  K_t^\alpha(A_tL_t)^{1-\alpha}.
% \]
% This upper bound is independent of the developer's choice of $B$.
% For a given locally integrable $r$, the discount factor
% $D_t=\exp(-\int_0^t r_s\dd s)$ has a positive minimum on $[0,T]$.
% Discounted operating revenue is bounded, whereas the cost of research
% tends to infinity with $\int_0^T M_t\dd t$. The developer's truncated
% supremum is therefore finite, and net payoff tends to minus infinity as
% total research expenditure diverges. This conclusion holds for every
% $0<\eta<1$, but does not prove that an infinite-horizon optimum exists.

\Needspace{7\baselineskip}
\begin{proof}[Proof of Proposition~\ref{cond:rewrite-uncapped-complements-limits}]
\phantomsection
\label{proof:rewrite-uncapped-complements-limits}
All auxiliary variables below are confined to this proof. Write
$x=X/(AL)$, $k=K/(AL)$, $c=C/(AL)$, $y=Y/(AL)$, $u=U/(AL)$,
$m=M/(AL)$, $d=\rho-n>0$, and $\nu=\delta+n+\gamma$.
No convergence of a growth rate or allocation ratio is assumed.

\textit{Step 1. Static optimality gives a limiting production function.}
Put $\varphi=(\sigma-1)/\sigma<0$ and
$f(x)=[\omega_L+\omega_Xx^\varphi]^{(1-\alpha)/\varphi}$, so $y=k^\alpha f(x)$.
The monopoly condition~\eqref{eq:rewrite-monopoly-foc} requires
$e_X=(1-s_X)/\sigma+\alpha s_X<1$. Thus
\begin{align}
 s_X>s_\infty&:=\frac{1-\sigma}{1-\alpha\sigma},&
 0<x<x_\infty&:=
 \left[\frac{\omega_Ls_\infty}{\omega_X(1-s_\infty)}\right]^{1/\varphi}.
 \label{eq:rewrite-uncapped-complements-x-limit}
\end{align}
Define $f_\infty=f(x_\infty)>0$.
In particular, $y\leq f_\infty k^\alpha$ and $0<u\leq x_\infty/B$.

The monopoly condition can be written
\[
 k^\alpha v(x)=1/B,\qquad
 v(x)=\frac{(1-\alpha)s_Xf(x)(1-e_X)}x.
\]
On $(0,x_\infty)$, $v$ is strictly decreasing, from infinity to zero.
Indeed, with a prime here denoting differentiation with respect to $\log x$,
$e_X'=(\alpha-\sigma^{-1})\varphi s_X(1-s_X)>0$ and
$-\mathrm d\log v/\mathrm d\log x=e_X+e_X'/(1-e_X)>\alpha s_X$.
Consequently $x$ increases with $B$ and $k$, while
\[
 \frac{\partial\log y}{\partial\log k}
 =\alpha+\frac{\alpha(1-\alpha)s_X}{e_X+e_X'/(1-e_X)}<1.
\]
Thus $y/k$ decreases with $k$ and increases with $B$.
For $k$ in any compact interval bounded away from zero, as $B\to\infty$,
\[
 x\to x_\infty,\qquad y\to f_\infty k^\alpha,\qquad u\to0
\]
uniformly. These are properties of the monopoly allocation at given $k,B$,
not assumptions about the equilibrium path. They also imply that, for any
prescribed finite return, sufficiently small $k$ and sufficiently large $B$
make $r$ exceed that return, uniformly over those small values of $k$.

\textit{Step 2. Household and developer optimality bound consumption.}
Proposition~\ref{prop:rewrite-uncapped-complements-bounds} bounds $k$ and
$y$ above by finite constants $k_{\max},y_{\max}$.
The household Euler equation and $r\geq-\delta$ give
$\dot c/c\geq-(\delta+\rho+\gamma)$.
Integrating the resource constraint over $[t,t+1]$ therefore gives
\[
 c_t\int_0^1 e^{-(\delta+\rho+\gamma)v}\dd v
 \leq\int_t^{t+1}c_s\dd s\leq k_{\max}+y_{\max}.
\]
Hence $c$ has a finite uniform upper bound $c_{\max}$.

Write $D_{t,s}=\exp(-\int_t^s r_v\dd v)$.
Euler implies $\int_t^\infty D_{t,s}C_s\dd s=C_t/d$.
Integrating the household budget~\eqref{eq:rewrite-household-budget}
and using its TVC~\eqref{eq:rewrite-household-tvc} shows that
the discounted integral of labor income plus developer net profit converges.
Since labor income is nonnegative, the discounted integral of developer net
profit has a limit that is either finite or minus infinity. The latter is impossible:
the developer can stop research and obtain nonnegative operating profits.
By continuation optimality, the developer's continuation value
$V_t=\int_t^\infty D_{t,s}\Pi_s\dd s$ is therefore finite and nonnegative.
It follows that
\begin{equation}
 \frac{C_t}{d}
 =K_t+V_t+\int_t^\infty D_{t,s}w_sL_s\dd s,\qquad
 c_t\geq dk_t,\qquad
 \int_t^\infty D_{t,s}w_sL_s\dd s\leq\frac{C_t}{d}.
 \label{eq:rewrite-complements-household-wealth-bound}
\end{equation}
Here $V_t$ is an auxiliary continuation value, not an additional traded
asset or a change to the household budget.

\textit{Step 3. Discounted research spending vanishes relative to effective labor.}
Fix a positive rate $a$. Step 1 permits a small $k_a>0$ and a sufficiently
late date such that $r-(n+\gamma)\geq a$ whenever $k\leq k_a$.
On $[k_a,k_{\max}]$, the wage equation~\eqref{eq:rewrite-wage} and the
static monopoly allocation give a uniform positive lower bound $\ell_a$ on $w/A$.
Since $r\geq-\delta$, there is a finite constant $H_a>0$ such that
\[
 a\leq r-(n+\gamma)+H_a w/A
\]
at every sufficiently late date. Multiplication by $D_{t,s}A_sL_s$,
integration, and~\eqref{eq:rewrite-complements-household-wealth-bound} give
\begin{equation}
 \int_t^\infty D_{t,s}A_sL_s\dd s
 \leq\frac{A_tL_t+H_aC_t/d}{a}
 \leq J A_tL_t
 \label{eq:rewrite-complements-effective-labor-pv-bound}
\end{equation}
for a fixed finite $J$. The omitted terminal contribution from discounted
effective labor is nonpositive; no discount-rate limit is used.

The research condition~\eqref{eq:rewrite-research-foc} and costate
equation~\eqref{eq:rewrite-costate} imply
\begin{equation}
 \frac{\mathrm d(qB)}{\mathrm dt}
 =rqB-U+\frac{1-\eta}{\eta}M.
 \label{eq:rewrite-complements-shadow-stock}
\end{equation}
Since $B$ is nondecreasing, $U_s\leq x_\infty A_sL_s/B_t$ for $s\geq t$.
Integrating~\eqref{eq:rewrite-complements-shadow-stock} first over finite
intervals and then using the developer TVC~\eqref{eq:rewrite-developer-tvc} gives
\begin{equation}
 q_tB_t+\frac{1-\eta}{\eta}
       \int_t^\infty D_{t,s}M_s\dd s
 =\int_t^\infty D_{t,s}U_s\dd s
 \leq\frac{x_\infty J}{B_t}A_tL_t.
 \label{eq:rewrite-complements-research-pv-bound}
\end{equation}
The right side is finite by~\eqref{eq:rewrite-complements-effective-labor-pv-bound};
nonnegativity then ensures that the research integral is finite as well.
Its value divided by $A_tL_t$ tends to zero. This argument requires
$0<\eta<1$, but neither $\eta\leq1/2$ nor $\eta<\alpha$.

\textit{Step 4. Capital cannot repeatedly approach zero.}
Let $F(k)=f_\infty k^\alpha-\nu k$.
Choose $\varepsilon>0$ small enough that $F$ is positive and increasing on
$(0,\varepsilon]$. Step 1 also permits a sufficiently late date and
$a>0$ such that $\dot c/c=r-\rho-\gamma\geq a$ whenever $k\leq\varepsilon$.
For any $j\leq\varepsilon$, the equilibrium path cannot enter
$\{k\leq j,\ c\geq F(j)\}$: there $\dot k\leq F(k)-c\leq0$,
whereas $c$ continues to grow at least at rate $a$, contradicting its upper bound.

If $k$ stayed below $\varepsilon$ forever, consumption would likewise
exceed its upper bound. Hence the path eventually reaches $\varepsilon$.
Suppose arbitrarily late crossings from $\varepsilon$ to $\varepsilon/2$
were possible. Let $t_0$ be the last crossing of $\varepsilon$ before the
first subsequent crossing $t_1$ of $\varepsilon/2$.
Throughout this interval $k\in[\varepsilon/2,\varepsilon]$ and $c$ increases.
Equation~\eqref{eq:rewrite-complements-household-wealth-bound} gives
$c_{t_0}\geq d\varepsilon$, so
\[
 t_1-t_0\leq T_\varepsilon
 :=a^{-1}\log\!\left(\frac{c_{\max}}{d\varepsilon}\right)<\infty.
\]
At $t_1$, the preceding invariant-region argument requires
$c_{t_1}<F(\varepsilon/2)$. Thus $c_s<F(\varepsilon/2)$ throughout the crossing.
The uniform convergence in Step 1, applied on the crossing interval, yields
numbers $\epsilon_{t_0}\to0$ such that
\[
 \dot k_s
 =y_s-u_s-\nu k_s-c_s-m_s
 \geq-\epsilon_{t_0}-m_s.
\]
On this interval $r$ is bounded above by a fixed positive $R_\varepsilon$.
Since $n+\gamma>0$, equation~\eqref{eq:rewrite-complements-research-pv-bound}
then implies
\[
 \frac{\varepsilon}{2}
 \leq T_\varepsilon\epsilon_{t_0}+\int_{t_0}^{t_1}m_s\dd s
 \leq T_\varepsilon\epsilon_{t_0}
 +e^{R_\varepsilon T_\varepsilon}
   \frac{\eta x_\infty J}{(1-\eta)B_{t_0}}
 \longrightarrow0,
\]
a contradiction. Therefore $k$ is eventually bounded away from zero,
and so is $c\geq dk$. In particular, $r$ is now bounded above.

\textit{Step 5. Research per unit of effective labor vanishes and Ramsey dynamics determine the limit.}
Differentiating the research condition and using the costate equation gives
\begin{equation}
 (1-\eta)g_M
 =r-\frac{\eta\chi UB^{\eta-1}}{M^{1-\eta}}\leq r.
 \label{eq:rewrite-complements-research-growth-bound}
\end{equation}
Choose a positive $R$ above all eventual values of $r$ and put $H=R/(1-\eta)$.
Then $M_s\geq M_t e^{-H(t-s)}$ on $[t-1,t]$ at sufficiently late dates.
Consequently
\[
 \frac1{A_{t-1}L_{t-1}}
 \int_{t-1}^t D_{t-1,s}M_s\dd s
 \geq e^{n+\gamma-R-H}m_t.
\]
The left side tends to zero by~\eqref{eq:rewrite-complements-research-pv-bound}.
Hence $m\to0$, and Step 1 already gives $u\to0$.

The normalized equilibrium equations therefore approach, uniformly on
the compact positive region containing the path, the Ramsey system
\begin{equation}
 \dot k=F(k)-c,\qquad
 \dot c=c\,[\alpha f_\infty k^{\alpha-1}-\delta-\rho-\gamma].
 \label{eq:rewrite-complements-limiting-ramsey}
\end{equation}
For clarity, the following argument verifies convergence rather than
inferring it solely from convergence of the equations.
This limiting system has one positive stationary point,
\[
 k_\infty=
 \left[\frac{\alpha f_\infty}{\rho+\gamma+\delta}\right]^{1/(1-\alpha)},
 \qquad c_\infty=F(k_\infty)>0.
\]
Its Jacobian has negative determinant, so it is a saddle.
The standard Ramsey phase portrait is discussed by \citet{groth2013ramsey}.
There is no nonconstant complete orbit confined to a compact subset of
the positive quadrant. To see this, the multiplier $1/c^2$ gives
\[
 \frac{\partial}{\partial k}\frac{F(k)-c}{c^2}
 +\frac{\partial}{\partial c}
   \frac{\alpha f_\infty k^{\alpha-1}-\delta-\rho-\gamma}{c}
 =\frac{\rho-n}{c^2}>0.
\]
The Bendixson--Dulac argument excludes both periodic orbits and homoclinic
loops. The planar limit-set theorem, the unique stationary point, and its
saddle character then exclude a nonconstant complete orbit confined to
such a compact set: its forward and backward limits would otherwise
require a periodic orbit or a connection back to that same saddle.

For any sequence of dates tending to infinity, time translations of our
bounded positive path have a subsequence converging on every bounded
time interval to a complete orbit of
\eqref{eq:rewrite-complements-limiting-ramsey}. This follows from bounded
derivatives, uniform convergence of the vector fields, and the integral
form of the differential equations. The preceding paragraph forces that
orbit to be the stationary point. Thus every accumulation point is
$(k_\infty,c_\infty)$, proving convergence of the original path.

It follows that $y\to y_\infty=f_\infty k_\infty^\alpha>0$ and
\begin{equation}
 r\to\rho+\gamma,\qquad
 \frac KY\to\frac{\alpha}{\rho+\gamma+\delta},\qquad
 \frac CY\to1-\frac{\alpha(\delta+n+\gamma)}{\rho+\gamma+\delta}>0.
 \label{eq:rewrite-complements-euler-limit}
\end{equation}
Positivity follows because the numerator of the final expression is
$\rho-n+(1-\alpha)(\delta+n+\gamma)>0$.
This proves~\eqref{eq:rewrite-uncapped-complements-allocation-ratios}.
Since $u,m\to0$ and $y\to y_\infty>0$, it follows that
$U/Y=u/y\to0$ and $M/Y=m/y\to0$, proving
\eqref{eq:rewrite-uncapped-complements-compute-shares}.
It also gives $\dot k\to0$ and $g_C\to n+\gamma$.

\textit{Step 6. Output and wage growth converge as well.}
A convergent level need not have a convergent derivative.
The bound on $g_M$ from Step 5 and the RSI technology give
\[
 B_t^{1-\eta}
 \geq(1-\eta)\chi M_t^\eta\int_0^1e^{-\eta H v}\dd v.
\]
Thus $g_B=\chi M^\eta/B^{1-\eta}$ is bounded above.
At $x_\infty$, the marginal-revenue function $v$ defined in Step 1 has
a simple zero: $v(x_\infty)=0$ and $v'(x_\infty)<0$, where the prime
now denotes differentiation with respect to $x$.
Differentiating $k^\alpha v(x)=1/B$ yields
\[
 \dot x=-\frac{v(x)}{v'(x)}
 \left(g_B+\alpha\frac{\dot k}{k}\right)\longrightarrow0.
\]
Hence $\dot y\to0$ and $\dot s_X\to0$.
Since $y_\infty>0$ and $1-s_\infty>0$, $g_{Y/N}\to\gamma$ and,
by~\eqref{eq:rewrite-wage}, $g_w\to\gamma$.
The limiting income shares follow by substituting $s_\infty$ into
\eqref{eq:rewrite-wage} and~\eqref{eq:rewrite-ai-price}.
The result is conditional on equilibrium existence and $B\to\infty$,
not an equilibrium-existence theorem or an assumption of balanced growth.
\end{proof}

\Needspace{7\baselineskip}
\begin{proof}[Proof of Proposition~\ref{prop:rewrite-uncapped-unit-bgp}]
\phantomsection
\label{proof:rewrite-uncapped-unit-bgp}
Use the definitions in Appendix~\ref{app:rewrite-uncapped-unit}.
\textit{The reference BGP.}
First, all candidate ratios are positive. The research denominator
$\rho-n+\eta g_Y^*$ is positive. Moreover,
\[
 i^*=\frac{\alpha(g_Y^*+\delta)}{\rho-n+g_Y^*+\delta}<\alpha,
 \qquad
 m^*<\frac{u^*\eta}{1-\eta}.
\]
Since $\Delta=(1-\alpha)(1-\eta-\omega_X)>0$,
$\omega_X<1-\eta$ and $\omega_X<1$. Consequently,
\[
 u^*+m^*<\frac{u^*}{1-\eta}<(1-\alpha)\omega_X<1-\alpha,
 \qquad c^*=1-u^*-i^*-m^*>0.
\]
Thus positive consumption follows from the stated conditions; it is not
an additional assumption.

For given $A_0,N_0>0$, choose the reference initial stocks as
\begin{align}
 B_0^*&=\left[
 A_0N_0(k^*)^{\alpha/[(1-\alpha)\omega_L]}(u^*)^{\omega_X/\omega_L}
 m^*\left(\frac{\chi}{g_B^*}\right)^{1/\eta}
 \right]^{\eta\omega_L/(1-\eta-\omega_X)},
 \label{eq:rewrite-uncapped-unit-initial-b}\\
 Y_0^*&=A_0N_0(k^*)^{\alpha/[(1-\alpha)\omega_L]}
 (u^*B_0^*)^{\omega_X/\omega_L},
 \qquad K_0^*=k^*Y_0^*.
 \label{eq:rewrite-uncapped-unit-initial-y-k}
\end{align}
The exponent is positive. These choices solve simultaneously the initial
production equation and
$g_B^*B_0^*=\chi(B_0^*m^*Y_0^*)^\eta$.
Set $U_0^*=u^*Y_0^*$, $M_0^*=m^*Y_0^*$, $C_0^*=c^*Y_0^*$,
$X_0^*=B_0^*U_0^*$, and
\begin{equation}
 q_0^*=\frac{M_0^*}{\eta g_B^*B_0^*}.
 \label{eq:rewrite-uncapped-unit-initial-q}
\end{equation}
Let $Y_t^*=Y_0^*e^{g_Y^*t}$, $B_t^*=B_0^*e^{g_B^*t}$,
$K_t^*=k^*Y_t^*$, $C_t^*=c^*Y_t^*$, $U_t^*=u^*Y_t^*$,
$M_t^*=m^*Y_t^*$, and $q_t^*=q_0^*e^{(g_Y^*-g_B^*)t}$.
Use the exogenous $A_t,N_t$ and set $L_t=N_t$ and $X_t^*=B_t^*U_t^*$.
The static equilibrium equations determine all remaining quantities and prices.

The two growth identities
\eqref{eq:rewrite-uncapped-unit-production-growth}
and~\eqref{eq:rewrite-uncapped-unit-research-growth} ensure that the production
equation and the RSI technology hold at every date once they hold initially.
The static monopoly condition gives $p_XX=(1-\alpha)\omega_XY$ and $U=u^*Y$.
Equation~\eqref{eq:rewrite-uncapped-unit-initial-q} gives
$M=\eta q\dot B$, which is the research first-order condition
\eqref{eq:rewrite-eq-research} with $\psi=1$.
The definitions of $k^*$, $r^*$, and $c^*$ verify the capital-return,
Euler, and resource equations. The costate equation reduces to
\[
 g_Y^*-g_B^*
 =r^*-\frac{u^*\eta g_B^*}{m^*}-\eta g_B^*.
\]
This holds by equation~\eqref{eq:rewrite-uncapped-unit-research-share}
and $(1-\eta)g_B^*=\eta g_Y^*$.
All allocations are positive and finite at every finite date.

Both transversality conditions are explicit:
\begin{align}
 e^{-\rho t}\frac{N_tK_t^*}{C_t^*}
 &=\frac{N_0k^*}{c^*}e^{-(\rho-n)t}\longrightarrow0,
 \label{eq:rewrite-uncapped-unit-household-tvc}\\
 e^{-r^*t}q_t^*B_t^*
 &=q_0^*B_0^*e^{-(\rho-n)t}\longrightarrow0.
 \label{eq:rewrite-uncapped-unit-developer-tvc}
\end{align}
Discounted developer profit is a constant multiple of
$e^{-(\rho-n)t}$. Discounted household utility is
$N_0e^{-(\rho-n)t}[\log(C_0^*/N_0)+(g_Y^*-n)t]$.
Both improper objective integrals are finite.

It remains to establish global optimality, rather than only the dated
necessary conditions. At given paths of $K,A,L$, static optimization gives
\begin{equation}
 \begin{split}
 \pi_t(B)={}&(1-\alpha)\omega_X[1-(1-\alpha)\omega_X]
 \bigl[K_t^\alpha(A_tL_t)^{(1-\alpha)\omega_L}\bigr]^{1/[1-(1-\alpha)\omega_X]}\\
 &\times\bigl[(1-\alpha)^2\omega_X^2B\bigr]^{(1-\alpha)\omega_X/[1-(1-\alpha)\omega_X]}.
 \end{split}
 \label{eq:rewrite-uncapped-unit-profit}
\end{equation}
For this verification, define the invertible state transformation and its
current-value costate by
\begin{equation}
 \nu=\frac{1-\eta}{\eta},\qquad
 S=B^\nu,\qquad p=\frac{q}{\nu B^{\nu-1}}.
 \label{eq:rewrite-uncapped-unit-transformed-state}
\end{equation}
The RSI technology becomes
\[
 \dot S=\nu\chi S^{1-\eta}M^\eta,
\]
which is jointly concave in $S,M$ for every $0<\eta<1$.
At given $K,A,L$, equation~\eqref{eq:rewrite-uncapped-unit-profit} is a
positive time-dependent coefficient times
\[
 S^{\frac{(1-\alpha)\omega_X}
 {\nu[1-(1-\alpha)\omega_X]}}.
\]
Since $\eta+\omega_X<1$ implies $(1-\alpha)\omega_X<1-\eta$, this
exponent lies strictly between zero and one. Operating profit is therefore
concave in $S$. With $p>0$, the transformed Hamiltonian
$\pi_t(S^{1/\nu})-M+p_t\nu\chi S^{1-\eta}M^\eta$ is jointly concave.
The research condition~\eqref{eq:rewrite-research-foc} and costate
equation~\eqref{eq:rewrite-costate}, with $\psi=1$, give its control
first-order condition and costate equation by the chain rule.

Write $D_t=\exp(-\int_0^t r_s\dd s)$ and
$J_T=\int_0^T D_t[\pi_t(B_t)-M_t]\dd t$.
For any admissible alternative with the same initial $B$, concavity,
the research condition, and the costate equation give, after integration
by parts,
\begin{equation}
 \widetilde J_T-J_T
 \leq-D_Tp_T(\widetilde S_T-S_T)
 \leq D_Tp_TS_T
 =\frac{D_Tq_TB_T}{\nu}\longrightarrow0.
 \label{eq:rewrite-uncapped-unit-developer-verification}
\end{equation}
The last limit follows from the developer TVC. Since the transformation
preserves feasible research paths and their payoffs, no admissible
alternative has a greater limiting payoff, or a greater upper limit of
truncated payoff, than this finite-valued candidate. The same verification
applies to nearby candidates satisfying the uncapped equations and the
developer TVC; it does not require exact balanced growth.

For the household, set $\Lambda_t=e^{-\rho t}N_t/C_t$.
Euler implies $\dot\Lambda=-r\Lambda$.
Concavity of log utility and the linear budget give, for any feasible
alternative with the same $K_0$ and the same given prices and distributions,
\[
 \int_0^T e^{-\rho t}N_t
 \left[\log(\widetilde C_t/N_t)-\log(C_t/N_t)\right]\dd t
 \leq-\Lambda_T(\widetilde K_T-K_T)
 \leq\Lambda_TK_T\longrightarrow0.
\]
This establishes global household optimality. Competitive final-good firms
face a concave constant-returns technology, and
Lemma~\ref{lem:rewrite-monopoly-choice} establishes global static monopoly
optimality. The constructed path satisfies Definition~\ref{def:rewrite-equilibrium}
with the uncapped RSI technology~\eqref{eq:rewrite-uncapped-unit-law}.

\medskip
\textit{Local convergence from nearby initial stocks.}
\phantomsection
\label{proof:rewrite-uncapped-unit-local}
Use the ordering $\xi=(\xi_K,\xi_B,\xi_C,\xi_q)'$ in
\eqref{eq:rewrite-uncapped-unit-deviations}, and let $e_K,e_B,e_C,e_q$
be the corresponding unit row vectors. Define the row vectors
\begin{align}
 y_\xi&=\left(\frac{\alpha}{1-(1-\alpha)\omega_X},
 \frac{(1-\alpha)\omega_X}{1-(1-\alpha)\omega_X},0,0\right),
 \label{eq:rewrite-uncapped-unit-output-row}\\
 m_\xi&=\frac{\eta e_B+e_q}{1-\eta},
 \label{eq:rewrite-uncapped-unit-research-row}\\
 b_\xi&=g_B^*[(\eta-1)e_B+\eta m_\xi].
 \label{eq:rewrite-uncapped-unit-efficiency-row}
\end{align}
The reduced production equation gives the log deviation of output as
$y_\xi\xi$. The research FOC implies
$M=(q\chi\eta B^\eta)^{1/(1-\eta)}$, so the log deviation of $M$ is
$m_\xi\xi$. Differentiating $g_B=\chi B^{\eta-1}M^\eta$ at the BGP
gives the level deviation $b_\xi\xi$.

The four rows of $J^*$, derived from capital accumulation, AI research,
the household Euler equation, and the developer costate equation, are
\begin{align}
 J_K^*&=\frac{1-u^*}{k^*}y_\xi-(g_Y^*+\delta)e_K
 -\frac{c^*}{k^*}e_C-\frac{m^*}{k^*}m_\xi,
 \label{eq:rewrite-uncapped-unit-jacobian-k}\\
 J_B^*&=b_\xi,
 \label{eq:rewrite-uncapped-unit-jacobian-b}\\
 J_C^*&=(r^*+\delta)(y_\xi-e_K),
 \label{eq:rewrite-uncapped-unit-jacobian-c}\\
 J_q^*&=(r^*+\delta)(y_\xi-e_K)
 -(\rho-n+\eta g_Y^*)(y_\xi-e_q-e_B)-\eta b_\xi.
 \label{eq:rewrite-uncapped-unit-jacobian-q}
\end{align}
The last row uses $U/(qB)=\rho-n+\eta g_Y^*$ at the BGP.
All entries depend on growth rates and ratios, which are independent of
$\chi$. No terminal restrictions were used to derive this Jacobian.

I next establish the two local properties
\eqref{eq:rewrite-uncapped-unit-local-spectrum}
and~\eqref{eq:rewrite-uncapped-unit-local-projection} from the parameter
restrictions, rather than impose them. Retain $\nu$ from
\eqref{eq:rewrite-uncapped-unit-transformed-state} and, for this linearization,
define
\begin{align*}
 a_K&=\frac{\alpha}{1-(1-\alpha)\omega_X},&
 a_T&=\frac{(1-\alpha)\omega_X\eta}
 {[1-(1-\alpha)\omega_X](1-\eta)},\\
 d&=\rho-n,&P&=r^*+\delta,&h&=\rho-n+\eta g_Y^*.
\end{align*}
All six quantities are positive, and the restriction $\eta+\omega_X<1$ gives
\begin{equation}
 1-a_K-a_T
 =\frac{(1-\alpha)(1-\eta-\omega_X)}
 {[1-(1-\alpha)\omega_X](1-\eta)}>0.
 \label{eq:rewrite-uncapped-unit-local-gap}
\end{equation}
Make the invertible linear change of deviations
\[
 \widehat\xi=(\xi_K,\nu\xi_B,\xi_C,\xi_q+(1-\nu)\xi_B)'.
\]
This changes neither the model nor the eigenvalues. Applying it to
\eqref{eq:rewrite-uncapped-unit-jacobian-k}--\eqref{eq:rewrite-uncapped-unit-jacobian-q}
gives a similar matrix
\begin{equation}
 \widehat J=
 \begin{pmatrix}
  \mathsf A&\mathsf B&-\mathsf C&-\mathsf V\\
  0&0&0&\mathsf D\\
  -\mathsf G&\mathsf H&0&0\\
  -\mathsf E&\mathsf F&0&d
 \end{pmatrix},
 \label{eq:rewrite-uncapped-unit-local-similarity}
\end{equation}
where the following letters abbreviate matrix coefficients only:
\begin{align*}
 \mathsf A&=d+(1-\alpha)\omega_X P,&
 \mathsf B&=\frac{(1-u^*)a_T-m^*}{k^*},&
 \mathsf C&=\frac{c^*}{k^*},\\
 \mathsf V&=\frac{m^*}{k^*(1-\eta)},&
 \mathsf D&=\frac{\eta g_Y^*}{1-\eta},&
 \mathsf G&=P(1-a_K),\\
 \mathsf H&=Pa_T,&
 \mathsf E&=\mathsf G+ha_K,&
 \mathsf F&=\mathsf H+h(1-a_T).
\end{align*}
All coefficients other than $\mathsf B$ are positive by their definitions.
The BGP allocation
ratios imply the identities
\begin{align}
 \mathsf L&\equiv\mathsf B-\frac{\mathsf Vd}{\mathsf D}
 =\frac{(1-\alpha)\omega_X\eta}{(1-\eta)k^*}>0,
 \label{eq:rewrite-uncapped-unit-local-l}\\
 \mathsf G\mathsf F-\mathsf H\mathsf E
 &=hP(1-a_K-a_T)>0.
 \label{eq:rewrite-uncapped-unit-local-minor}
\end{align}
Set $\mathsf N=\mathsf C\mathsf G+\mathsf V\mathsf E>0$ and
$\mathsf Q=\mathsf C\mathsf H+\mathsf V\mathsf F>0$.
Expanding the determinant gives the characteristic polynomial
\begin{equation}
 \begin{split}
 p(\ell)&=\det(\ell I-\widehat J)\\
 &=(\ell^2-\mathsf A\ell-\mathsf N)
 (\ell^2-d\ell-\mathsf D\mathsf F)
 -\mathsf D\mathsf E(\mathsf Q-\mathsf L\ell).
 \end{split}
 \label{eq:rewrite-uncapped-unit-local-polynomial}
\end{equation}
In particular,
\[
 p(0)=\mathsf D\mathsf C hP(1-a_K-a_T)>0.
\]
Let $\mu=(d-\sqrt{d^2+4\mathsf D\mathsf F})/2<0$, the negative
root of the second quadratic factor. Then
$p(\mu)=-\mathsf D\mathsf E(\mathsf Q-\mathsf L\mu)<0$.
Since $p(\ell)\to+\infty$ as $\ell\to-\infty$, there is one
negative root below $\mu$ and another between $\mu$ and zero.
These are the only negative roots: the cubic coefficient of $p$ is
$-(\mathsf A+d)<0$, its linear coefficient is
$\mathsf A\mathsf D\mathsf F+d\mathsf N+\mathsf L\mathsf D\mathsf E>0$,
and its constant coefficient is positive. Thus the coefficient signs of
$p(-x)$, in descending powers, are $(+,+,\text{arbitrary},-,+)$.
Descartes' rule of signs gives exactly two sign changes, regardless of
the quadratic coefficient. The two negative roots are therefore real,
distinct, and simple. Write them as $\ell_1<\mu<\ell_2<0$.
The remaining two roots have positive product and positive sum, because
$p(0)>0$ and the sum of all four roots is $\mathsf A+d>0$.
If real, both are positive; if complex, their real parts are positive.
This proves \eqref{eq:rewrite-uncapped-unit-local-spectrum}.

To establish the state projection, write a stable eigenvector in transformed
coordinates as $(x,y,z,u)'$. Its second component cannot be zero:
the second, fourth, and first rows would successively imply $u=x=z=0$.
Normalize $y=1$. The second and fourth rows then give
\[
 x(\ell)=\frac{\mathsf D\mathsf F+d\ell-\ell^2}{\mathsf D\mathsf E}.
\]
The ordering of the roots implies $x(\ell_1)<0<x(\ell_2)$.
Hence the two projected stable eigenvectors $(x(\ell_1),1)'$ and
$(x(\ell_2),1)'$ are linearly independent. Returning to $(\xi_K,\xi_B)$
only rescales their second component by $1/\nu>0$, so
\eqref{eq:rewrite-uncapped-unit-local-projection} follows as well.

The exact normalized vector field is continuously differentiable near
$\xi=0$. The spectral property just proved and the stable-manifold
theorem supply a two-dimensional local manifold
whose trajectories remain near zero and converge exponentially.
The derivative of its projection onto $(\xi_K,\xi_B)$ is $V_{KB}$.
The nonsingularity just proved and the
inverse-function theorem therefore make this manifold locally the
graph of a continuously differentiable function selecting $(\xi_C,\xi_q)$
from the two initial stock deviations. Since exponentiation is a local
diffeomorphism, the same holds for $(C_0,q_0)$ as functions of $(K_0,B_0)$.
More explicitly, for some constants $a,H>0$, trajectories on a sufficiently
small stable manifold satisfy
\[
 \|\xi(t)\|\leq H e^{-at}\|\xi(0)\|,\qquad t\geq0.
\]
For any sufficiently small open neighborhood $\mathcal V$ of zero,
choose a ball contained in $\mathcal V$. Continuity of the graph selecting
the two jump variables permits an open neighborhood $\mathcal U$ of the
initial stock deviations for which $H\|\xi(0)\|$ is smaller than that
ball's radius. Every initial stock pair in $\mathcal U$ therefore has a
trajectory staying in $\mathcal V$ for all $t\geq0$ and converging to zero.
Exponentiation maps $\mathcal U$ to an open neighborhood of
$(K_0^*,B_0^*)$. Shrinking these neighborhoods preserves positivity,
feasibility, and existence for every $t\geq0$. Local uniqueness follows
because every trajectory remaining in this local region and converging
to zero belongs to the same stable manifold, which is a graph over stocks.
This does not assert convergence from arbitrary initial $C_0,q_0$ or
exclude equilibrium trajectories that leave the region.

The statement also permits nearby $A_0,N_0$. Equations
\eqref{eq:rewrite-uncapped-unit-initial-b}
and~\eqref{eq:rewrite-uncapped-unit-initial-y-k} make the reference
$B_0^*,K_0^*$ smooth functions of these two positive initial values.
After normalization by the corresponding BGP, the exact vector field
depends only on the parameters and BGP ratios, not on $A_0,N_0$.
The set of $(A_0,N_0,K_0,B_0)$ whose stock deviations belong to
$\mathcal U$ is therefore open and contains the reference initial condition.
From each such initial condition, the construction converges to the BGP
associated with its own $A_0,N_0$, not necessarily to the same levels as
the original reference path.

Since $f(\xi(t))\to0$, growth rates converge to their BGP values, not only
normalized levels. The static equations give the same convergence for
the interest rate, wage growth, and allocation ratios. The unit-elastic
labor income share remains $(1-\alpha)\omega_L$ at every date.

Exponential convergence implies that $r_t-r^*$ is integrable. Thus
$D_t/e^{-r^*t}$ tends to a finite positive constant, and both TVC expressions
equal their BGP counterparts in
\eqref{eq:rewrite-uncapped-unit-household-tvc}
and~\eqref{eq:rewrite-uncapped-unit-developer-tvc} times factors tending
to finite positive limits. Both TVCs hold. Output and all expenditure
flows are bounded multiples of $e^{g_Y^*t}$, while log consumption per
person is a linear function of time plus a bounded deviation. Both
objective integrals therefore remain finite. The global concavity
verification established above
applies at these new price and allocation paths, proving agent optimality.
Static optimality and market clearing then establish equilibrium, not
merely a solution of the necessary conditions.
\end{proof}


\Needspace{6\baselineskip}
\begin{proof}[Proof of Proposition~\ref{prop:rewrite-research-scale}]
\label{proof:rewrite-research-scale}
Write $\theta=(1-\alpha)/\alpha$. For a research policy $M$, discounted
net profit on the fixed horizon is
\[
 J_T(M)=\int_0^T
 \exp\left\{-\int_0^t r_s\dd s\right\}
 [\pi(B_t)-M_t]\dd t.
\]
Integrability of $r$ bounds the discount factor above and away from zero.
Since $\sigma>1$, optimized operating profit satisfies
\[
 \pi(B)\sim\kappa_XKB^\theta
 \qquad\text{as }B\to\infty,
\]
for a constant $\kappa_X>0$, uniformly along the given positive continuous
paths of $K,A,$ and $L$ on $[0,T]$.

Fix $\Delta\in(0,T)$, choose $M=\mathcal M$ on $[0,\Delta]$, and choose
$M=0$ afterward. The uncapped RSI technology gives
\[
 B_{\mathcal M}(t)^{1-\eta}=B_0^{1-\eta}
 +(1-\eta)\chi\mathcal M^\eta\min\{t,\Delta\}.
\]
Since optimized operating profit is nonnegative, a lower bound on discounted
net profit counts only operating profits earned after research stops, while
subtracting all research expenditure. For all sufficiently large $\mathcal M$,
\begin{equation}
 J_T(\mathcal M)\geq c_\pi\mathcal M^{
 \eta(1-\alpha)/[\alpha(1-\eta)]}
 -c_M\mathcal M,
 \label{eq:rewrite-research-burst-payoff}
\end{equation}
where $c_\pi,c_M>0$ are independent of $\mathcal M$.
The exponent exceeds one because $\eta>\alpha$,
so $J_T(\mathcal M)\to+\infty$ as $\mathcal M\to\infty$.
Each finite $\mathcal M$ keeps AI efficiency and discounted net profit
finite on $[0,T]$; it is the sequence of payoffs that is unbounded.
Continuing with a candidate equilibrium's research expenditure after $T$
cannot reduce subsequent profits once the burst leaves AI efficiency above
the candidate's level.
\end{proof}

% BEGIN PRESERVED BROADER RESEARCH-SCALE PROOF (2026-09-15)
% The active proposition now states only the eta > alpha finite-horizon result.
% Keep the earlier eta < alpha bound, boundary case, finite-bound existence,
% and infinite-horizon deviation argument available without compiling them.
% \paragraph{Research expenditure and finite-horizon value.}
%
% Besides the two cases in Proposition~\ref{prop:rewrite-research-scale},
% the proof gives auxiliary finite-horizon bounds for the same given paths.
% At $\eta=\alpha$, the leading benefit and cost terms of the research-burst
% construction below have the same order, so comparing exponents alone does
% not settle the result. With a fixed finite upper bound, the truncated problem
% has a finite value and attains its maximum for every $0<\eta<1$.
% Global verification of a candidate research policy is addressed separately
% in Lemma~\ref{lem:rewrite-developer-verification}.
%
% \Needspace{6\baselineskip}
% \begin{proof}[Proof of Proposition~\ref{prop:rewrite-research-scale}]
% \label{proof:rewrite-research-scale}
% \emph{Part (i).} I first consider the uncapped economy. Define
% \[
%  \theta=\frac{1-\alpha}{\alpha},\qquad
%  S=\int_0^T M_t\dd t.
% \]
% Here $S$ is total research expenditure over the fixed horizon $[0,T]$.
% For an admissible research policy, let its discounted net payoff be
% \[
%  J_T(M)=\int_0^T
%  \exp\left\{-\int_0^t r_s\dd s\right\}
%  [\pi(B_t)-M_t]\dd t.
% \]
% For every admissible policy, \(B\) is nondecreasing. The RSI technology implies
% \[
%  \frac{\dd}{\dd t}B_t^{1-\eta}
%  =(1-\eta)\chi M_t^\eta.
% \]
% Integration and H\"older's inequality give
% \begin{align*}
%  B_t^{1-\eta}
%  &=B_0^{1-\eta}+(1-\eta)\chi\int_0^t M_s^\eta\dd s,\\
%  \int_0^T M_t^\eta\dd t
%  &\leq T^{1-\eta}S^\eta.
% \end{align*}
% Because $0<\eta<1$, spreading a given expenditure evenly over $[0,T]$
% maximizes the integral of $M^\eta$. Hence there is a constant $C_B>0$,
% independent of the research policy,
% such that
% \begin{equation}
%  \max_{0\leq t\leq T}B_t
%  \leq C_B\left(1+S^{\eta/(1-\eta)}\right),
%  \label{eq:rewrite-capability-expenditure-bound}
% \end{equation}
% The same bound covers
% arbitrarily concentrated controls because, on a support of length \(d\leq T\),
% \(\int M_t^\eta\dd t\leq d^{1-\eta}S^\eta\).
%
% Let \(\overline K=\max_{t\leq T}K_t\) and
% \(\overline{\mathcal L}=\max_{t\leq T}A_tL_t\). The weighted power-mean
% inequality gives \(Z\leq AL+X\). The final-good firm's demand condition therefore
% implies
% \[
%  p_XX\leq(1-\alpha)\overline K^\alpha
%  \left(\overline{\mathcal L}^{1-\alpha}+X^{1-\alpha}\right).
% \]
% \Needspace{6\baselineskip}
% Write the right-hand side as $a+bX^{1-\alpha}$, with $a,b>0$ determined by
% the given paths and production parameters. The first-order condition for
% maximizing this revenue bound net of $X/B$ is
% $b(1-\alpha)X^{-\alpha}=1/B$, and its maximized value is
% \[
%  \max_{X\geq0}\left\{a+bX^{1-\alpha}-\frac{X}{B}\right\}
%  =a+\alpha b[b(1-\alpha)B]^\theta.
% \]
% There is therefore a constant $C_\pi>0$, independent of the research policy,
% such that
% \begin{equation}
%  \pi(B)\leq C_\pi(1+B^\theta).
%  \label{eq:rewrite-profit-upper-bound}
% \end{equation}
% Combining the AI-efficiency and profit bounds shows that operating profit grows
% no faster than a constant plus a term proportional to
% \begin{equation}
%  S^{\eta(1-\alpha)/[\alpha(1-\eta)]},
%  \label{eq:rewrite-research-benefit-order}
% \end{equation}
% where the exponent is the product of $\eta/(1-\eta)$, from expenditure to
% AI efficiency, and $(1-\alpha)/\alpha$, from AI efficiency to profit. This is an
% upper bound on the order of growth, not an exact profit function.
%
% On a finite horizon, integrability of $r$ bounds the discount factor above
% and away from zero. Let $\underline D>0$ be its minimum on $[0,T]$.
% Combining
% \eqref{eq:rewrite-capability-expenditure-bound} and
% \eqref{eq:rewrite-profit-upper-bound} yields, for constants $C_0,C_1>0$
% independent of the research policy,
% \begin{equation}
%  J_T(M)\leq C_0+C_1S^p-\underline D S,
%  \qquad
%  p=\frac{\eta(1-\alpha)}{\alpha(1-\eta)}.
%  \label{eq:rewrite-truncated-value-bound}
% \end{equation}
% The exponent satisfies $p<1$ exactly when $\eta<\alpha$. In this case, the
% benefit grows at most sublinearly in expenditure while its cost grows linearly.
% The bound tends to minus infinity as $S\to\infty$, so the value is finite and
% the objective is coercive.
%
% \emph{Part (ii).} An upper bound with $p>1$ does not by itself prove unboundedness. For the
% converse, I construct research policies with unbounded payoffs. Suppose
% \(\sigma>1\). As \(B\to\infty\), optimized
% operating profit satisfies
% \[
%  \pi(B)\sim\kappa_XKB^\theta
% \]
% for a constant \(\kappa_X>0\), uniformly when \(K,A,\) and \(L\) remain in
% compact subsets of the positive orthant. Fix \(\Delta\in(0,T)\), choose
% \(M=\mathcal M\) on \([0,\Delta]\), and choose \(M=0\) afterward. Then
% \[
%  B_{\mathcal M}(t)^{1-\eta}=B_0^{1-\eta}
%  +(1-\eta)\chi\mathcal M^\eta\min\{t,\Delta\}.
% \]
% Choosing \(X\) optimally at every date gives
% \begin{equation}
%  J_T(\mathcal M)=c_\pi\mathcal M^{
%  \eta(1-\alpha)/[\alpha(1-\eta)]}
%  -c_M\mathcal M
%  +o\!\left(\mathcal M^{
%  \eta(1-\alpha)/[\alpha(1-\eta)]}\right),
%  \label{eq:rewrite-research-burst-payoff}
% \end{equation}
% with \(c_\pi,c_M>0\). The exponent exceeds one when \(\eta>\alpha\), so this
% family makes the truncated value unbounded above. At equality, the exponent is
% one and the sign of the leading coefficient determines this family's result.
%
% To obtain the infinite-horizon nonexistence conclusion, suppose instead that
% an uncapped equilibrium has a developer policy $\widehat M$, an associated
% efficiency path $\widehat B$, and a finite discounted value $\widehat V_0$.
% On every finite interval of this candidate equilibrium, $K,A,L$ are positive
% and continuous, and $r=\alpha Y/K-\delta$ is continuous by the static
% production and monopoly solution. The finite-horizon argument therefore
% applies to these paths, without requiring a BGP or a bounded interest rate
% as $t\to\infty$.
% Hold its paths of $K,A,L,$ and $r$ fixed, as required by the developer's
% problem~\eqref{eq:rewrite-developer-reduced}. On $[0,T]$, use the burst policy
% above; after $T$, resume $\widehat M$. For a sufficiently large burst,
% $B_{\mathcal M}(T)>\widehat B(T)$. Under the uncapped law, the difference
% $B_{\mathcal M}(t)^{1-\eta}-\widehat B(t)^{1-\eta}$ is then constant and
% positive for $t\geq T$, since the two policies use the same research spending.
% Operating profit is increasing in $B$ by
% equation~\eqref{eq:rewrite-profit-envelope}, so subsequent net profit is no
% lower than under the candidate policy. Let $\widehat J_T$ be the candidate's
% discounted net profit on $[0,T]$ and let $J_\infty(\mathcal M)$ denote the
% payoff from the deviating policy, including its continuation. Then
% \[
%  J_\infty(\mathcal M)
%  \geq J_T(\mathcal M)+\widehat V_0-\widehat J_T
%  \longrightarrow+\infty.
% \]
% The continuation payoff is well defined as a finite value or $+\infty$:
% its difference from the candidate's continuation payoff is the integral of
% a nonnegative operating-profit gain. Each deviation has finite research
% expenditure on every finite interval and hence finite $B$ there by
% \eqref{eq:rewrite-capability-expenditure-bound}. Every finite-valued candidate
% can therefore be improved, contradicting developer optimality. The argument
% excludes a finite-valued optimum; it does not treat an infinite payoff as an
% attained firm value or assume that a trajectory reaches a singularity.
%
% For a finite upper bound, the invariance result above confines \(B\) to the compact
% interval \([B_0,\overline B]\). Optimized operating profit
% is continuous in \(B\) and uniformly bounded on the fixed horizon. Hence
% \(J_T(M)\leq C_0-\underline D S\), so every maximizing sequence has bounded
% total expenditure. Set \(v=M^\eta\) and \(p=1/\eta>1\). The control cost is
% \(v^p\), while the state equation is affine in \(v\):
% \[
%  \dot B=\chi B^\eta\left(1-\frac{B}{\overline B}\right)v.
% \]
% A maximizing sequence is bounded in the reflexive space \(L^p[0,T]\). A weakly
% convergent subsequence and the associated uniformly convergent state paths
% preserve feasibility. The operating-profit integral is continuous under
% uniform state convergence, while \(\int v^p\) is weakly lower semicontinuous.
% The objective is therefore weakly upper semicontinuous and attains its maximum.
% \end{proof}
%
% \Needspace{4\baselineskip}
% The condition $\eta<\alpha$ prevents arbitrarily large
% research expenditure from being optimal on a fixed horizon in the uncapped
% economy. It is a sufficient restriction, not a necessary one for every value
% of $\sigma$, and finite-horizon existence with a finite upper bound does not
% require $\eta<\alpha$. Neither the bound for $\eta<\alpha$ nor
% finite-horizon developer existence with a finite upper bound establishes an
% equilibrium: the paths of $K,A,L,$ and $r$ were taken as given. In contrast,
% for $\sigma>1$ and $\eta>\alpha$, the profitable-deviation argument rules out
% an uncapped equilibrium with a finite-valued developer optimum at any such
% candidate paths.
% The quantitative analysis uses $\eta=0.20<\alpha=0.33$, for which the
% expenditure exponent is approximately $0.508$; the value of $\eta$ is
% illustrative rather than estimated.
% END PRESERVED BROADER RESEARCH-SCALE PROOF

\Needspace{9\baselineskip}
\begin{proposition}[An AI bottleneck under complementarity]
\label{prop:rewrite-low-cap-complements}
Suppose $0<\sigma<1$ and $0<B_0<\overline B<\mathcal B(\sigma)$.
Any equilibrium satisfies
\begin{gather*}
 \limsup_{t\to\infty}\frac1t
 \log\left(\frac{Y_t/N_t}{Y_0/N_0}\right)
 \leq\mathcal R(\overline B)-\rho<\gamma,\\
 \frac{X}{AL}\to0,\qquad s_X\to1,\qquad \frac{wL}{Y}\to0.
\end{gather*}
\end{proposition}

\Needspace{4\baselineskip}
\begin{proof}[Proof of Proposition~\ref{prop:rewrite-low-cap-complements}]
\phantomsection
\label{proof:rewrite-low-cap-complements}
Consider any equilibrium with $0<\sigma<1$ and
$0<B_0<\overline B<\mathcal B(\sigma)$. The increasing share map in
equation~\eqref{eq:rewrite-frontier-share-map}, established in Step 1 of the
proof of Proposition~\ref{prop:rewrite-equilibrium-regimes},
and $B_t\leq\overline B$ imply
\begin{equation}
 r_t\leq\mathcal R(B_t)\leq\mathcal R(\overline B)<\rho+\gamma,
 \qquad
 \frac{Y_t}{K_t}\leq\frac{\mathcal R(\overline B)+\delta}{\alpha}.
 \label{eq:rewrite-low-cap-return-bound}
\end{equation}
The developer's discounted continuation net profit is nonnegative: it can
stop research and earn nonnegative operating profit. Wage income is also
nonnegative. Integrating the household budget
\eqref{eq:rewrite-household-budget}, using its transversality condition
\eqref{eq:rewrite-household-tvc} and the Euler equation
\eqref{eq:rewrite-euler}, therefore gives the first inequality below;
integrating the Euler equation with the return bound gives the second:
\begin{equation}
 C_t\geq(\rho-n)K_t,
 \qquad
 \frac{C_t}{N_t}\leq\frac{C_0}{N_0}
 e^{[\mathcal R(\overline B)-\rho]t}.
 \label{eq:rewrite-low-cap-consumption-bound}
\end{equation}
Combining these inequalities yields
\begin{equation}
 \frac{Y_t}{N_t}\leq
 \frac{\mathcal R(\overline B)+\delta}{\alpha(\rho-n)}
 \frac{C_0}{N_0}e^{[\mathcal R(\overline B)-\rho]t},
 \qquad
 \limsup_{t\to\infty}\frac1t
 \log\left(\frac{Y_t/N_t}{Y_0/N_0}\right)
 \leq\mathcal R(\overline B)-\rho<\gamma.
 \label{eq:rewrite-low-cap-growth-bound}
\end{equation}
Consequently $Y/(AL)\to0$. Since equation~\eqref{eq:rewrite-ai-output-limit}
gives $X/Y\leq\overline B(1-\alpha)^2$ under complementarity,
$X/(AL)\to0$: AI production services become scarce relative to effective
labor. Equation~\eqref{eq:rewrite-sx} then implies $s_X\to1$ and hence
$wL/Y\to0$.
The growth bound concerns the long-run average, not convergence of the
instantaneous growth rate. Neither convergence of $C/Y$ nor
$B_t\to\overline B$ is assumed. These are necessary properties of any
equilibrium in this region, not a proof of equilibrium existence.
\end{proof}
% The previous Section 5.3 and its proposition are commented in the source.
% Keep the complete proof file available without compiling an orphan proof.
% \input{sections_rewrite/appendix_uncapped_substitutes_proof}

\section{Numerical algorithm and equilibrium audit}
\label{app:rewrite-algorithm}

\subsection{Economic intuition and initial conditions}

The BVP selects an equilibrium trajectory, not an initial-value solution
from arbitrary consumption and shadow prices. Capital and AI efficiency
are predetermined; consumption and the developer's shadow value adjust
to put the economy on the stable continuation of its long-run regime.
Satisfying the static conditions alone does not select that continuation.

For each $\sigma$, I first solve the existing-AI economy with efficiency
fixed at $B_0$ and research unavailable. A scalar root search in $\log K_0$
solves Equation~\eqref{eq:rewrite-rsi-initial-ratio} using the production
and monopoly equations. The resource constraint~\eqref{eq:rewrite-resource},
with $M=0$ and $\dot K=(n+\gamma)K$, gives pre-event consumption.
Homogeneity extends these values to a BGP:
$K,Y,C,U,X$ grow at $n+\gamma$, wages grow at $\gamma$, and
$B,p_X,r$ and income shares are constant. Since
$r-(n+\gamma)=\rho-n>0$, pre-event developer profit has finite present
value and the household's transversality expression decays.

At unexpected RSI activation, only the states are inherited. The
positive-$\chi$ problem selects consumption, research and the shadow
value anew. A separate audit checks continuity of production quantities
and prices, verifies the pre-event Euler, resource and monopoly conditions,
and records the endogenous jumps in consumption and research.

\subsection{Boundary-value solution and reported outcomes}

I use the collocation solver \texttt{scipy.integrate.solve\_bvp} with
computational coordinates
\[
 \left(\log K,\ \log\frac{B}{\overline B-B},\ \log C,\ \log q\right).
\]
Subtracting the analytical terminal aggregate growth rate times $t$
removes their linear trends without changing model variables or timing.
At each trial state the code solves the static monopoly problem and
recovers research from its first-order condition. The left boundary fixes
$K_0,B_0$; two right-boundary projections remove unstable components of
the normalized equilibrium dynamics near the long-run limit, using
Appendix~\ref{app:rewrite-finite-existence}.

For each of the eight parameter configurations, I start near its
analytical limiting point and move the initial stocks to their prescribed
values in 32 continuation stages. Neither $\omega_X=0$ nor $\chi=0$
is used as a degenerate BVP seed. At unit elasticity the static block
uses its exact Cobb--Douglas limit; near one, \texttt{log1p} and
\texttt{expm1} avoid cancellation. Under complementarity, marginal
revenue is evaluated directly and the static root is polished when
needed near zero marginal revenue; the first-order condition is unchanged.

The values $\chi=7.5$ and $\chi=1.5$ are supplied directly: there is no
outer calibration search. Research shares and price declines are outputs.
Annual research expenditure relative to output is measured as
$\int_0^1 M_t\,\dd t/\int_0^1Y_t\,\dd t$, not $M_0/Y_0$.
Gauss--Legendre quadrature with 64 and 128 nodes must agree within
$10^{-8}$ in logs. The annual ratio must also change by less than
$2\times10^{-5}$ in logs across the horizon extensions.

\subsection{Equilibrium checks}

I extend each computational horizon twice by 500 years and tighten the
collocation tolerance from $2\times10^{-6}$ to $10^{-8}$ and $10^{-9}$.
The final boundary tolerance is $10^{-11}$. Initial meshes contain
221, 401 and 601 points and adapt as needed. The displayed window lies
inside each solved horizon; no spline is extrapolated.

The original equations are reconstructed independently with five-point
differences at 1,001 dates and steps of 0.003 and 0.001 years. An additional
801-date check over the first ten years uses steps of 0.0003 and 0.0001.
Admission requires dynamic residuals below $10^{-6}$, first-order
residuals below $10^{-9}$, final-horizon changes below $2\times10^{-5}$
in detrended logs, and terminal-coordinate gaps below $10^{-4}$.
These are the existing numerical tolerances, not relaxed economic conditions.

The audit checks feasibility and developer global optimality. When
global profit concavity in research depth fails, it checks the Hamiltonian
support condition in Lemma~\ref{lem:rewrite-developer-verification},
using counterfactual-state minimization and an analytical bound beyond
the grid. The two grids contain 81/101 and 321/241 date/state points,
supplemented by dense first-decade checks when needed. The analytical
support bound and the asymptotic rate $n-\rho<0$ of both transversality
expressions support the infinite-horizon continuation. These are
numerically verified equilibrium approximations, not interval-arithmetic
existence certificates for arbitrary initial stocks.

\begin{table}[H]
\centering
\caption{Numerical accuracy of the illustrative RSI comparisons}
\label{tab:rewrite-rsi-illustrative-accuracy}
\small
\begin{tabular}{rrrrrr}
\toprule
$\chi$ & $\sigma$ & Final horizon & Mesh points & Dynamic residual & Horizon change\\
\midrule
7.5 & 0.90 & 5,428.0 & 1,193 & $1.18\times10^{-9}$ & $3.06\times10^{-8}$\\
7.5 & 1.00 & 4,944.5 & 1,290 & $1.38\times10^{-9}$ & $3.46\times10^{-8}$\\
7.5 & 1.10 & 4,828.8 & 1,326 & $1.66\times10^{-9}$ & $3.73\times10^{-8}$\\
7.5 & 1.50 & 3,013.0 & 1,586 & $1.69\times10^{-9}$ & $1.56\times10^{-8}$\\
1.5 & 0.90 & 6,047.0 & 841 & $1.02\times10^{-9}$ & $4.30\times10^{-8}$\\
1.5 & 1.00 & 5,563.5 & 893 & $1.11\times10^{-9}$ & $3.80\times10^{-8}$\\
1.5 & 1.10 & 5,447.8 & 911 & $1.19\times10^{-9}$ & $4.08\times10^{-8}$\\
1.5 & 1.50 & 3,632.0 & 1,104 & $1.23\times10^{-9}$ & $1.80\times10^{-8}$\\
\bottomrule
\end{tabular}
\par\smallskip
\begin{minipage}{0.98\textwidth}\footnotesize
Horizons are measured in years. Dynamic residuals are the largest independently reconstructed errors across
the full-horizon and dense first-decade checks. Horizon change compares the
last two solutions in detrended logarithmic coordinates over years 0--500.
Error bounds are rounded upward. All eight paths pass the equation,
event-continuity, developer-optimality and long-run-continuation checks.
The $\sigma=1.50$ cases use the global Hamiltonian-support test and its
analytical continuation bound; the other cases pass the concavity test.
\end{minipage}
\end{table}


\subsection{Replication files}

The code, simulation data, and replication instructions are available at
\url{https://github.com/oliverpardo1979/ai-growth-and-labor}.
The \href{https://github.com/oliverpardo1979/ai-growth-and-labor/blob/main/REPLICATION.md}
{replication guide} explains how to install the dependencies, check the
published results, and recompute the eight trajectories. From the repository
root, run:
\begin{quote}
\small
\begin{verbatim}
python scripts/simulate_rewrite_illustrative_rsi.py
\end{verbatim}
\end{quote}
This command solves and audits all four elasticities in both scenarios.
Use \texttt{--chi 7.5} or \texttt{--chi 1.5} for one scenario;
\texttt{--sigma 1.5}, for example, solves one elasticity. With all four
checkpoints available, \texttt{--finish} repeats the admission checks
and exports the figures. The full replication entry point is
\texttt{scripts/reproduce\_rewrite\_results.py}.

The \href{https://github.com/oliverpardo1979/ai-growth-and-labor/tree/main/numerical_rewrite/rsi_chi_7_5}
{higher-productivity results folder} and
\href{https://github.com/oliverpardo1979/ai-growth-and-labor/tree/main/numerical_rewrite/rsi_chi_1_5}
{lower-productivity results folder} contain the scenario inputs, admitted
CSV, outcome summaries, equilibrium checks and figure metadata.
Checkpoint and CSV hashes bind the audits to the plotted data.
No comparison is exported unless all four scenarios and the activation
continuity checks pass.

The repository includes the two comparisons reported here. Earlier
experiments are preserved separately and are not needed for replication.

